% Môn: Vật lý
% Lớp: 10
% Chương: Chương III. Động lực học
% Bài: L10C3 Bài 21. Mômen lực. Cân bằng của vật rắn
% Số câu: 162
% App đọc file này trực tiếp — sửa rồi Commit trên GitHub, bấm Đồng bộ GitHub .tex

% L10C3 Bài 21. Mômen lực. Cân bằng của vật rắn

% ===== Câu 1 =====
% ID: AIQ_L10C3_FULL_1297
% Mức: NB
\dangbt{Nhận biết mômen lực và cánh tay đòn}
\begin{ex}
% ID: AIQ_L10C3_FULL_1297
% Mức: NB
Đại lượng nào quyết định tác dụng làm quay của lực đối với một trục?
\choice
{\True Tích của lực và cánh tay đòn}
{Chỉ độ lớn của lực}
{Chỉ khối lượng vật}
{Thời gian tác dụng}
\loigiai{
Mômen lực đối với trục là $M=Fd$, với d là cánh tay đòn. Chọn A.
}
\end{ex}

% ===== Câu 2 =====
% ID: AIQ_L10C3_FULL_1298
% Mức: NB
\begin{ex}
% ID: AIQ_L10C3_FULL_1298
% Mức: NB
Đại lượng nào quyết định tác dụng làm quay của lực đối với một trục?
\choice
{\True Tích của lực và cánh tay đòn}
{Chỉ độ lớn của lực}
{Chỉ khối lượng vật}
{Thời gian tác dụng}
\loigiai{
Mômen lực đối với trục là $M=Fd$, với d là cánh tay đòn. Chọn A.
}
\end{ex}

% ===== Câu 3 =====
% ID: AIQ_L10C3_FULL_1299
% Mức: NB
\begin{ex}
% ID: AIQ_L10C3_FULL_1299
% Mức: NB
Đại lượng nào quyết định tác dụng làm quay của lực đối với một trục?
\choice
{\True Tích của lực và cánh tay đòn}
{Chỉ độ lớn của lực}
{Chỉ khối lượng vật}
{Thời gian tác dụng}
\loigiai{
Mômen lực đối với trục là $M=Fd$, với d là cánh tay đòn. Chọn A.
}
\end{ex}

% ===== Câu 4 =====
% ID: AIQ_L10C3_FULL_1300
% Mức: TH
\begin{ex}
% ID: AIQ_L10C3_FULL_1300
% Mức: TH
Đại lượng nào quyết định tác dụng làm quay của lực đối với một trục?
\choice
{\True Tích của lực và cánh tay đòn}
{Chỉ độ lớn của lực}
{Chỉ khối lượng vật}
{Thời gian tác dụng}
\loigiai{
Mômen lực đối với trục là $M=Fd$, với d là cánh tay đòn. Chọn A.
}
\end{ex}

% ===== Câu 5 =====
% ID: AIQ_L10C3_FULL_1301
% Mức: TH
\begin{ex}
% ID: AIQ_L10C3_FULL_1301
% Mức: TH
Đại lượng nào quyết định tác dụng làm quay của lực đối với một trục?
\choice
{\True Tích của lực và cánh tay đòn}
{Chỉ độ lớn của lực}
{Chỉ khối lượng vật}
{Thời gian tác dụng}
\loigiai{
Mômen lực đối với trục là $M=Fd$, với d là cánh tay đòn. Chọn A.
}
\end{ex}

% ===== Câu 6 =====
% ID: AIQ_L10C3_FULL_1302
% Mức: TH
\begin{ex}
% ID: AIQ_L10C3_FULL_1302
% Mức: TH
Đại lượng nào quyết định tác dụng làm quay của lực đối với một trục?
\choice
{\True Tích của lực và cánh tay đòn}
{Chỉ độ lớn của lực}
{Chỉ khối lượng vật}
{Thời gian tác dụng}
\loigiai{
Mômen lực đối với trục là $M=Fd$, với d là cánh tay đòn. Chọn A.
}
\end{ex}

% ===== Câu 7 =====
% ID: AIQ_L10C3_FULL_1303
% Mức: TH
\begin{ex}
% ID: AIQ_L10C3_FULL_1303
% Mức: TH
Đại lượng nào quyết định tác dụng làm quay của lực đối với một trục?
\choice
{\True Tích của lực và cánh tay đòn}
{Chỉ độ lớn của lực}
{Chỉ khối lượng vật}
{Thời gian tác dụng}
\loigiai{
Mômen lực đối với trục là $M=Fd$, với d là cánh tay đòn. Chọn A.
}
\end{ex}

% ===== Câu 8 =====
% ID: AIQ_L10C3_FULL_1304
% Mức: VD
\begin{ex}
% ID: AIQ_L10C3_FULL_1304
% Mức: VD
Đại lượng nào quyết định tác dụng làm quay của lực đối với một trục?
\choice
{\True Tích của lực và cánh tay đòn}
{Chỉ độ lớn của lực}
{Chỉ khối lượng vật}
{Thời gian tác dụng}
\loigiai{
Mômen lực đối với trục là $M=Fd$, với d là cánh tay đòn. Chọn A.
}
\end{ex}

% ===== Câu 9 =====
% ID: AIQ_L10C3_FULL_1305
% Mức: VD
\begin{ex}
% ID: AIQ_L10C3_FULL_1305
% Mức: VD
Đại lượng nào quyết định tác dụng làm quay của lực đối với một trục?
\choice
{\True Tích của lực và cánh tay đòn}
{Chỉ độ lớn của lực}
{Chỉ khối lượng vật}
{Thời gian tác dụng}
\loigiai{
Mômen lực đối với trục là $M=Fd$, với d là cánh tay đòn. Chọn A.
}
\end{ex}

% ===== Câu 10 =====
% ID: AIQ_L10C3_FULL_1306
% Mức: VDC
\begin{ex}
% ID: AIQ_L10C3_FULL_1306
% Mức: VDC
Đại lượng nào quyết định tác dụng làm quay của lực đối với một trục?
\choice
{\True Tích của lực và cánh tay đòn}
{Chỉ độ lớn của lực}
{Chỉ khối lượng vật}
{Thời gian tác dụng}
\loigiai{
Mômen lực đối với trục là $M=Fd$, với d là cánh tay đòn. Chọn A.
}
\end{ex}

% ===== Câu 11 =====
% ID: AIQ_L10C3_FULL_1307
% Mức: NB
\begin{ex}
% ID: AIQ_L10C3_FULL_1307
% Mức: NB
Đại lượng nào quyết định tác dụng làm quay của lực đối với một trục?
\choiceTF
{\True Lực cản phụ thuộc chuyển động tương đối giữa vật và chất lưu.}
{Lực cản luôn cùng chiều chuyển động của vật.}
{\True Khi hợp lực bằng 0, vật có gia tốc bằng 0.}
{Có thể bỏ qua đơn vị của đại lượng trong kết quả vật lí.}
\loigiai{
\begin{itemchoice}
\item (Đúng) Lực cản phụ thuộc chuyển động tương đối giữa vật và chất lưu. (Vì): Mômen lực đối với trục là $M=Fd$, với d là cánh tay đòn. b) Sai vì giá trị hoặc chiều nêu ra không phù hợp. c) Đúng theo định luật II Newton. d) Sai vì kết quả vật lí phải kèm đơn vị thích hợp
\item (Sai) Lực cản luôn cùng chiều chuyển động của vật.
\item (Đúng) Khi hợp lực bằng 0, vật có gia tốc bằng 0.
\item (Sai) Có thể bỏ qua đơn vị của đại lượng trong kết quả vật lí.
\end{itemchoice}
}
\end{ex}

% ===== Câu 12 =====
% ID: AIQ_L10C3_FULL_1308
% Mức: NB
\begin{ex}
% ID: AIQ_L10C3_FULL_1308
% Mức: NB
Đại lượng nào quyết định tác dụng làm quay của lực đối với một trục?
\choiceTF
{\True Lực cản phụ thuộc chuyển động tương đối giữa vật và chất lưu.}
{Lực cản luôn cùng chiều chuyển động của vật.}
{\True Khi hợp lực bằng 0, vật có gia tốc bằng 0.}
{Có thể bỏ qua đơn vị của đại lượng trong kết quả vật lí.}
\loigiai{
\begin{itemchoice}
\item (Đúng) Lực cản phụ thuộc chuyển động tương đối giữa vật và chất lưu. (Vì): Mômen lực đối với trục là $M=Fd$, với d là cánh tay đòn. b) Sai vì giá trị hoặc chiều nêu ra không phù hợp. c) Đúng theo định luật II Newton. d) Sai vì kết quả vật lí phải kèm đơn vị thích hợp
\item (Sai) Lực cản luôn cùng chiều chuyển động của vật.
\item (Đúng) Khi hợp lực bằng 0, vật có gia tốc bằng 0.
\item (Sai) Có thể bỏ qua đơn vị của đại lượng trong kết quả vật lí.
\end{itemchoice}
}
\end{ex}

% ===== Câu 13 =====
% ID: AIQ_L10C3_FULL_1309
% Mức: TH
\begin{ex}
% ID: AIQ_L10C3_FULL_1309
% Mức: TH
Đại lượng nào quyết định tác dụng làm quay của lực đối với một trục?
\choiceTF
{\True Lực cản phụ thuộc chuyển động tương đối giữa vật và chất lưu.}
{Lực cản luôn cùng chiều chuyển động của vật.}
{\True Khi hợp lực bằng 0, vật có gia tốc bằng 0.}
{Có thể bỏ qua đơn vị của đại lượng trong kết quả vật lí.}
\loigiai{
\begin{itemchoice}
\item (Đúng) Lực cản phụ thuộc chuyển động tương đối giữa vật và chất lưu. (Vì): Mômen lực đối với trục là $M=Fd$, với d là cánh tay đòn. b) Sai vì giá trị hoặc chiều nêu ra không phù hợp. c) Đúng theo định luật II Newton. d) Sai vì kết quả vật lí phải kèm đơn vị thích hợp
\item (Sai) Lực cản luôn cùng chiều chuyển động của vật.
\item (Đúng) Khi hợp lực bằng 0, vật có gia tốc bằng 0.
\item (Sai) Có thể bỏ qua đơn vị của đại lượng trong kết quả vật lí.
\end{itemchoice}
}
\end{ex}

% ===== Câu 14 =====
% ID: AIQ_L10C3_FULL_1310
% Mức: TH
\begin{ex}
% ID: AIQ_L10C3_FULL_1310
% Mức: TH
Đại lượng nào quyết định tác dụng làm quay của lực đối với một trục?
\choiceTF
{\True Lực cản phụ thuộc chuyển động tương đối giữa vật và chất lưu.}
{Lực cản luôn cùng chiều chuyển động của vật.}
{\True Khi hợp lực bằng 0, vật có gia tốc bằng 0.}
{Có thể bỏ qua đơn vị của đại lượng trong kết quả vật lí.}
\loigiai{
\begin{itemchoice}
\item (Đúng) Lực cản phụ thuộc chuyển động tương đối giữa vật và chất lưu. (Vì): Mômen lực đối với trục là $M=Fd$, với d là cánh tay đòn. b) Sai vì giá trị hoặc chiều nêu ra không phù hợp. c) Đúng theo định luật II Newton. d) Sai vì kết quả vật lí phải kèm đơn vị thích hợp
\item (Sai) Lực cản luôn cùng chiều chuyển động của vật.
\item (Đúng) Khi hợp lực bằng 0, vật có gia tốc bằng 0.
\item (Sai) Có thể bỏ qua đơn vị của đại lượng trong kết quả vật lí.
\end{itemchoice}
}
\end{ex}

% ===== Câu 15 =====
% ID: AIQ_L10C3_FULL_1311
% Mức: VD
\begin{ex}
% ID: AIQ_L10C3_FULL_1311
% Mức: VD
Đại lượng nào quyết định tác dụng làm quay của lực đối với một trục?
\choiceTF
{\True Lực cản phụ thuộc chuyển động tương đối giữa vật và chất lưu.}
{Lực cản luôn cùng chiều chuyển động của vật.}
{\True Khi hợp lực bằng 0, vật có gia tốc bằng 0.}
{Có thể bỏ qua đơn vị của đại lượng trong kết quả vật lí.}
\loigiai{
\begin{itemchoice}
\item (Đúng) Lực cản phụ thuộc chuyển động tương đối giữa vật và chất lưu. (Vì): Mômen lực đối với trục là $M=Fd$, với d là cánh tay đòn. b) Sai vì giá trị hoặc chiều nêu ra không phù hợp. c) Đúng theo định luật II Newton. d) Sai vì kết quả vật lí phải kèm đơn vị thích hợp
\item (Sai) Lực cản luôn cùng chiều chuyển động của vật.
\item (Đúng) Khi hợp lực bằng 0, vật có gia tốc bằng 0.
\item (Sai) Có thể bỏ qua đơn vị của đại lượng trong kết quả vật lí.
\end{itemchoice}
}
\end{ex}

% ===== Câu 16 =====
% ID: AIQ_L10C3_FULL_1312
% Mức: TH
\begin{ex}
% ID: AIQ_L10C3_FULL_1312
% Mức: TH
Đại lượng nào quyết định tác dụng làm quay của lực đối với một trục?
\shortans{Tích}
\loigiai{
Mômen lực đối với trục là $M=Fd$, với d là cánh tay đòn. Đáp án trả lời ngắn không ghi đơn vị.
}
\end{ex}

% ===== Câu 17 =====
% ID: AIQ_L10C3_FULL_1313
% Mức: TH
\begin{ex}
% ID: AIQ_L10C3_FULL_1313
% Mức: TH
Đại lượng nào quyết định tác dụng làm quay của lực đối với một trục?
\shortans{Tích}
\loigiai{
Mômen lực đối với trục là $M=Fd$, với d là cánh tay đòn. Đáp án trả lời ngắn không ghi đơn vị.
}
\end{ex}

% ===== Câu 18 =====
% ID: AIQ_L10C3_FULL_1314
% Mức: VD
\begin{ex}
% ID: AIQ_L10C3_FULL_1314
% Mức: VD
Đại lượng nào quyết định tác dụng làm quay của lực đối với một trục?
\shortans{Tích}
\loigiai{
Mômen lực đối với trục là $M=Fd$, với d là cánh tay đòn. Đáp án trả lời ngắn không ghi đơn vị.
}
\end{ex}

% ===== Câu 19 =====
% ID: AIQ_L10C3_FULL_1315
% Mức: VD
\begin{ex}
% ID: AIQ_L10C3_FULL_1315
% Mức: VD
Đại lượng nào quyết định tác dụng làm quay của lực đối với một trục?
\shortans{Tích}
\loigiai{
Mômen lực đối với trục là $M=Fd$, với d là cánh tay đòn. Đáp án trả lời ngắn không ghi đơn vị.
}
\end{ex}

% ===== Câu 20 =====
% ID: AIQ_L10C3_FULL_1316
% Mức: VD
\begin{ex}
% ID: AIQ_L10C3_FULL_1316
% Mức: VD
Đại lượng nào quyết định tác dụng làm quay của lực đối với một trục?
\shortans{Tích}
\loigiai{
Mômen lực đối với trục là $M=Fd$, với d là cánh tay đòn. Đáp án trả lời ngắn không ghi đơn vị.
}
\end{ex}

% ===== Câu 21 =====
% ID: AIQ_L10C3_FULL_1317
% Mức: VDC
\begin{ex}
% ID: AIQ_L10C3_FULL_1317
% Mức: VDC
Đại lượng nào quyết định tác dụng làm quay của lực đối với một trục?
\shortans{Tích}
\loigiai{
Mômen lực đối với trục là $M=Fd$, với d là cánh tay đòn. Đáp án trả lời ngắn không ghi đơn vị.
}
\end{ex}

% ===== Câu 22 =====
% ID: AIQ_L10C3_FULL_1318
% Mức: TH
\begin{bt}
% ID: AIQ_L10C3_FULL_1318
% Mức: TH
Đại lượng nào quyết định tác dụng làm quay của lực đối với một trục?
\loigiai{
Phân tích dữ kiện, chọn định luật phù hợp. Mômen lực đối với trục là $M=Fd$, với d là cánh tay đòn.
}
\end{bt}

% ===== Câu 23 =====
% ID: AIQ_L10C3_FULL_1319
% Mức: TH
\begin{bt}
% ID: AIQ_L10C3_FULL_1319
% Mức: TH
Đại lượng nào quyết định tác dụng làm quay của lực đối với một trục?
\loigiai{
Phân tích dữ kiện, chọn định luật phù hợp. Mômen lực đối với trục là $M=Fd$, với d là cánh tay đòn.
}
\end{bt}

% ===== Câu 24 =====
% ID: AIQ_L10C3_FULL_1320
% Mức: VD
\begin{bt}
% ID: AIQ_L10C3_FULL_1320
% Mức: VD
Đại lượng nào quyết định tác dụng làm quay của lực đối với một trục?
\loigiai{
Phân tích dữ kiện, chọn định luật phù hợp. Mômen lực đối với trục là $M=Fd$, với d là cánh tay đòn.
}
\end{bt}

% ===== Câu 25 =====
% ID: AIQ_L10C3_FULL_1321
% Mức: VD
\begin{bt}
% ID: AIQ_L10C3_FULL_1321
% Mức: VD
Đại lượng nào quyết định tác dụng làm quay của lực đối với một trục?
\loigiai{
Phân tích dữ kiện, chọn định luật phù hợp. Mômen lực đối với trục là $M=Fd$, với d là cánh tay đòn.
}
\end{bt}

% ===== Câu 26 =====
% ID: AIQ_L10C3_FULL_1322
% Mức: VD
\begin{bt}
% ID: AIQ_L10C3_FULL_1322
% Mức: VD
Đại lượng nào quyết định tác dụng làm quay của lực đối với một trục?
\loigiai{
Phân tích dữ kiện, chọn định luật phù hợp. Mômen lực đối với trục là $M=Fd$, với d là cánh tay đòn.
}
\end{bt}

% ===== Câu 27 =====
% ID: AIQ_L10C3_FULL_1323
% Mức: VDC
\begin{bt}
% ID: AIQ_L10C3_FULL_1323
% Mức: VDC
Đại lượng nào quyết định tác dụng làm quay của lực đối với một trục?
\loigiai{
Phân tích dữ kiện, chọn định luật phù hợp. Mômen lực đối với trục là $M=Fd$, với d là cánh tay đòn.
}
\end{bt}

% ===== Câu 28 =====
% ID: AIQ_L10C3_FULL_1324
% Mức: NB
\dangbt{Tính mômen của lực đối với trục quay}
\begin{ex}
% ID: AIQ_L10C3_FULL_1324
% Mức: NB
Lực 13 $N$ có giá cách trục quay $0,4\,\text{m}$. Tính mômen lực.
\choice
{\True $5,2\,\text{N}$.m}
{$6,24\,\text{N}$.m}
{$4,16\,\text{N}$.m}
{$7,2\,\text{N}$.m}
\loigiai{
$M=Fd=13\cdot0,4=5,2$ N.m. Chọn A.
}
\end{ex}

% ===== Câu 29 =====
% ID: AIQ_L10C3_FULL_1325
% Mức: NB
\begin{ex}
% ID: AIQ_L10C3_FULL_1325
% Mức: NB
Lực 15 $N$ có giá cách trục quay $0,45\,\text{m}$. Tính mômen lực.
\choice
{$8,1\,\text{N}$.m}
{$5,4\,\text{N}$.m}
{$8,75\,\text{N}$.m}
{\True $6,75\,\text{N}$.m}
\loigiai{
$M=Fd=15\cdot0,45=6,75$ N.m. Chọn D.
}
\end{ex}

% ===== Câu 30 =====
% ID: AIQ_L10C3_FULL_1326
% Mức: NB
\begin{ex}
% ID: AIQ_L10C3_FULL_1326
% Mức: NB
Lực 17 $N$ có giá cách trục quay $0,2\,\text{m}$. Tính mômen lực.
\choice
{$2,72\,\text{N}$.m}
{$5,4\,\text{N}$.m}
{\True $3,4\,\text{N}$.m}
{$4,08\,\text{N}$.m}
\loigiai{
$M=Fd=17\cdot0,2=3,4$ N.m. Chọn C.
}
\end{ex}

% ===== Câu 31 =====
% ID: AIQ_L10C3_FULL_1327
% Mức: TH
\begin{ex}
% ID: AIQ_L10C3_FULL_1327
% Mức: TH
Lực 19 $N$ có giá cách trục quay $0,25\,\text{m}$. Tính mômen lực.
\choice
{$6,75\,\text{N}$.m}
{\True $4,75\,\text{N}$.m}
{$5,7\,\text{N}$.m}
{$3,8\,\text{N}$.m}
\loigiai{
$M=Fd=19\cdot0,25=4,75$ N.m. Chọn B.
}
\end{ex}

% ===== Câu 32 =====
% ID: AIQ_L10C3_FULL_1328
% Mức: TH
\begin{ex}
% ID: AIQ_L10C3_FULL_1328
% Mức: TH
Lực 13 $N$ có giá cách trục quay $0,3\,\text{m}$. Tính mômen lực.
\choice
{\True $3,9\,\text{N}$.m}
{$4,68\,\text{N}$.m}
{$3,12\,\text{N}$.m}
{$5,9\,\text{N}$.m}
\loigiai{
$M=Fd=13\cdot0,3=3,9$ N.m. Chọn A.
}
\end{ex}

% ===== Câu 33 =====
% ID: AIQ_L10C3_FULL_1329
% Mức: TH
\begin{ex}
% ID: AIQ_L10C3_FULL_1329
% Mức: TH
Lực 15 $N$ có giá cách trục quay $0,35\,\text{m}$. Tính mômen lực.
\choice
{$6,3\,\text{N}$.m}
{$4,2\,\text{N}$.m}
{$7,25\,\text{N}$.m}
{\True $5,25\,\text{N}$.m}
\loigiai{
$M=Fd=15\cdot0,35=5,25$ N.m. Chọn D.
}
\end{ex}

% ===== Câu 34 =====
% ID: AIQ_L10C3_FULL_1330
% Mức: TH
\begin{ex}
% ID: AIQ_L10C3_FULL_1330
% Mức: TH
Lực 17 $N$ có giá cách trục quay $0,4\,\text{m}$. Tính mômen lực.
\choice
{$5,44\,\text{N}$.m}
{$8,8\,\text{N}$.m}
{\True $6,8\,\text{N}$.m}
{$8,16\,\text{N}$.m}
\loigiai{
$M=Fd=17\cdot0,4=6,8$ N.m. Chọn C.
}
\end{ex}

% ===== Câu 35 =====
% ID: AIQ_L10C3_FULL_1331
% Mức: VD
\begin{ex}
% ID: AIQ_L10C3_FULL_1331
% Mức: VD
Lực 19 $N$ có giá cách trục quay $0,45\,\text{m}$. Tính mômen lực.
\choice
{$10,55\,\text{N}$.m}
{\True $8,55\,\text{N}$.m}
{$10,26\,\text{N}$.m}
{$6,84\,\text{N}$.m}
\loigiai{
$M=Fd=19\cdot0,45=8,55$ N.m. Chọn B.
}
\end{ex}

% ===== Câu 36 =====
% ID: AIQ_L10C3_FULL_1332
% Mức: VD
\begin{ex}
% ID: AIQ_L10C3_FULL_1332
% Mức: VD
Lực 13 $N$ có giá cách trục quay $0,2\,\text{m}$. Tính mômen lực.
\choice
{\True $2,6\,\text{N}$.m}
{$3,12\,\text{N}$.m}
{$2,08\,\text{N}$.m}
{$4,6\,\text{N}$.m}
\loigiai{
$M=Fd=13\cdot0,2=2,6$ N.m. Chọn A.
}
\end{ex}

% ===== Câu 37 =====
% ID: AIQ_L10C3_FULL_1333
% Mức: VDC
\begin{ex}
% ID: AIQ_L10C3_FULL_1333
% Mức: VDC
Lực 15 $N$ có giá cách trục quay $0,25\,\text{m}$. Tính mômen lực.
\choice
{$4,5\,\text{N}$.m}
{$3\,\text{N}$.m}
{$5,75\,\text{N}$.m}
{\True $3,75\,\text{N}$.m}
\loigiai{
$M=Fd=15\cdot0,25=3,75$ N.m. Chọn D.
}
\end{ex}

% ===== Câu 38 =====
% ID: AIQ_L10C3_FULL_1334
% Mức: NB
\begin{ex}
% ID: AIQ_L10C3_FULL_1334
% Mức: NB
Lực 17 $N$ có giá cách trục quay $0,3\,\text{m}$. Tính mômen lực.
\choiceTF
{\True Kết quả đúng là $5,1\,\text{N}$.m.}
{Kết quả bằng $10,2\,\text{N}$.m.}
{\True Khi hợp lực bằng 0, vật có gia tốc bằng 0.}
{Có thể bỏ qua đơn vị của đại lượng trong kết quả vật lí.}
\loigiai{
\begin{itemchoice}
\item (Đúng) Kết quả đúng là $5,1\,\text{N}$.m. (Vì): $M=Fd=17\cdot0,3=5,1$ N.m. b) Sai vì giá trị hoặc chiều nêu ra không phù hợp. c) Đúng theo định luật II Newton. d) Sai vì kết quả vật lí phải kèm đơn vị thích hợp
\item (Sai) Kết quả bằng $10,2\,\text{N}$.m.
\item (Đúng) Khi hợp lực bằng 0, vật có gia tốc bằng 0.
\item (Sai) Có thể bỏ qua đơn vị của đại lượng trong kết quả vật lí.
\end{itemchoice}
}
\end{ex}

% ===== Câu 39 =====
% ID: AIQ_L10C3_FULL_1335
% Mức: NB
\begin{ex}
% ID: AIQ_L10C3_FULL_1335
% Mức: NB
Lực 19 $N$ có giá cách trục quay $0,35\,\text{m}$. Tính mômen lực.
\choiceTF
{\True Kết quả đúng là $6,65\,\text{N}$.m.}
{Kết quả bằng $13,3\,\text{N}$.m.}
{\True Khi hợp lực bằng 0, vật có gia tốc bằng 0.}
{Có thể bỏ qua đơn vị của đại lượng trong kết quả vật lí.}
\loigiai{
\begin{itemchoice}
\item (Đúng) Kết quả đúng là $6,65\,\text{N}$.m. (Vì): $M=Fd=19\cdot0,35=6,65$ N.m. b) Sai vì giá trị hoặc chiều nêu ra không phù hợp. c) Đúng theo định luật II Newton. d) Sai vì kết quả vật lí phải kèm đơn vị thích hợp
\item (Sai) Kết quả bằng $13,3\,\text{N}$.m.
\item (Đúng) Khi hợp lực bằng 0, vật có gia tốc bằng 0.
\item (Sai) Có thể bỏ qua đơn vị của đại lượng trong kết quả vật lí.
\end{itemchoice}
}
\end{ex}

% ===== Câu 40 =====
% ID: AIQ_L10C3_FULL_1336
% Mức: TH
\begin{ex}
% ID: AIQ_L10C3_FULL_1336
% Mức: TH
Lực 13 $N$ có giá cách trục quay $0,4\,\text{m}$. Tính mômen lực.
\choiceTF
{\True Kết quả đúng là $5,2\,\text{N}$.m.}
{Kết quả bằng $10,4\,\text{N}$.m.}
{\True Khi hợp lực bằng 0, vật có gia tốc bằng 0.}
{Có thể bỏ qua đơn vị của đại lượng trong kết quả vật lí.}
\loigiai{
\begin{itemchoice}
\item (Đúng) Kết quả đúng là $5,2\,\text{N}$.m. (Vì): $M=Fd=13\cdot0,4=5,2$ N.m. b) Sai vì giá trị hoặc chiều nêu ra không phù hợp. c) Đúng theo định luật II Newton. d) Sai vì kết quả vật lí phải kèm đơn vị thích hợp
\item (Sai) Kết quả bằng $10,4\,\text{N}$.m.
\item (Đúng) Khi hợp lực bằng 0, vật có gia tốc bằng 0.
\item (Sai) Có thể bỏ qua đơn vị của đại lượng trong kết quả vật lí.
\end{itemchoice}
}
\end{ex}

% ===== Câu 41 =====
% ID: AIQ_L10C3_FULL_1337
% Mức: TH
\begin{ex}
% ID: AIQ_L10C3_FULL_1337
% Mức: TH
Lực 15 $N$ có giá cách trục quay $0,45\,\text{m}$. Tính mômen lực.
\choiceTF
{\True Kết quả đúng là $6,75\,\text{N}$.m.}
{Kết quả bằng $13,5\,\text{N}$.m.}
{\True Khi hợp lực bằng 0, vật có gia tốc bằng 0.}
{Có thể bỏ qua đơn vị của đại lượng trong kết quả vật lí.}
\loigiai{
\begin{itemchoice}
\item (Đúng) Kết quả đúng là $6,75\,\text{N}$.m. (Vì): $M=Fd=15\cdot0,45=6,75$ N.m. b) Sai vì giá trị hoặc chiều nêu ra không phù hợp. c) Đúng theo định luật II Newton. d) Sai vì kết quả vật lí phải kèm đơn vị thích hợp
\item (Sai) Kết quả bằng $13,5\,\text{N}$.m.
\item (Đúng) Khi hợp lực bằng 0, vật có gia tốc bằng 0.
\item (Sai) Có thể bỏ qua đơn vị của đại lượng trong kết quả vật lí.
\end{itemchoice}
}
\end{ex}

% ===== Câu 42 =====
% ID: AIQ_L10C3_FULL_1338
% Mức: VD
\begin{ex}
% ID: AIQ_L10C3_FULL_1338
% Mức: VD
Lực 17 $N$ có giá cách trục quay $0,2\,\text{m}$. Tính mômen lực.
\choiceTF
{\True Kết quả đúng là $3,4\,\text{N}$.m.}
{Kết quả bằng $6,8\,\text{N}$.m.}
{\True Khi hợp lực bằng 0, vật có gia tốc bằng 0.}
{Có thể bỏ qua đơn vị của đại lượng trong kết quả vật lí.}
\loigiai{
\begin{itemchoice}
\item (Đúng) Kết quả đúng là $3,4\,\text{N}$.m. (Vì): $M=Fd=17\cdot0,2=3,4$ N.m. b) Sai vì giá trị hoặc chiều nêu ra không phù hợp. c) Đúng theo định luật II Newton. d) Sai vì kết quả vật lí phải kèm đơn vị thích hợp
\item (Sai) Kết quả bằng $6,8\,\text{N}$.m.
\item (Đúng) Khi hợp lực bằng 0, vật có gia tốc bằng 0.
\item (Sai) Có thể bỏ qua đơn vị của đại lượng trong kết quả vật lí.
\end{itemchoice}
}
\end{ex}

% ===== Câu 43 =====
% ID: AIQ_L10C3_FULL_1339
% Mức: TH
\begin{ex}
% ID: AIQ_L10C3_FULL_1339
% Mức: TH
Lực 19 $N$ có giá cách trục quay $0,25\,\text{m}$. Tính mômen lực.
\shortans{4,75}
\loigiai{
$M=Fd=19\cdot0,25=4,75$ N.m. Đáp án trả lời ngắn không ghi đơn vị.
}
\end{ex}

% ===== Câu 44 =====
% ID: AIQ_L10C3_FULL_1340
% Mức: TH
\begin{ex}
% ID: AIQ_L10C3_FULL_1340
% Mức: TH
Lực 13 $N$ có giá cách trục quay $0,3\,\text{m}$. Tính mômen lực.
\shortans{3,9}
\loigiai{
$M=Fd=13\cdot0,3=3,9$ N.m. Đáp án trả lời ngắn không ghi đơn vị.
}
\end{ex}

% ===== Câu 45 =====
% ID: AIQ_L10C3_FULL_1341
% Mức: VD
\begin{ex}
% ID: AIQ_L10C3_FULL_1341
% Mức: VD
Lực 15 $N$ có giá cách trục quay $0,35\,\text{m}$. Tính mômen lực.
\shortans{5,25}
\loigiai{
$M=Fd=15\cdot0,35=5,25$ N.m. Đáp án trả lời ngắn không ghi đơn vị.
}
\end{ex}

% ===== Câu 46 =====
% ID: AIQ_L10C3_FULL_1342
% Mức: VD
\begin{ex}
% ID: AIQ_L10C3_FULL_1342
% Mức: VD
Lực 17 $N$ có giá cách trục quay $0,4\,\text{m}$. Tính mômen lực.
\shortans{6,8}
\loigiai{
$M=Fd=17\cdot0,4=6,8$ N.m. Đáp án trả lời ngắn không ghi đơn vị.
}
\end{ex}

% ===== Câu 47 =====
% ID: AIQ_L10C3_FULL_1343
% Mức: VD
\begin{ex}
% ID: AIQ_L10C3_FULL_1343
% Mức: VD
Lực 19 $N$ có giá cách trục quay $0,45\,\text{m}$. Tính mômen lực.
\shortans{8,55}
\loigiai{
$M=Fd=19\cdot0,45=8,55$ N.m. Đáp án trả lời ngắn không ghi đơn vị.
}
\end{ex}

% ===== Câu 48 =====
% ID: AIQ_L10C3_FULL_1344
% Mức: VDC
\begin{ex}
% ID: AIQ_L10C3_FULL_1344
% Mức: VDC
Lực 13 $N$ có giá cách trục quay $0,2\,\text{m}$. Tính mômen lực.
\shortans{2,6}
\loigiai{
$M=Fd=13\cdot0,2=2,6$ N.m. Đáp án trả lời ngắn không ghi đơn vị.
}
\end{ex}

% ===== Câu 49 =====
% ID: AIQ_L10C3_FULL_1345
% Mức: TH
\begin{bt}
% ID: AIQ_L10C3_FULL_1345
% Mức: TH
Lực 15 $N$ có giá cách trục quay $0,25\,\text{m}$. Tính mômen lực.
\loigiai{
Phân tích dữ kiện, chọn định luật phù hợp. $M=Fd=15\cdot0,25=3,75$ N.m.
}
\end{bt}

% ===== Câu 50 =====
% ID: AIQ_L10C3_FULL_1346
% Mức: TH
\begin{bt}
% ID: AIQ_L10C3_FULL_1346
% Mức: TH
Lực 17 $N$ có giá cách trục quay $0,3\,\text{m}$. Tính mômen lực.
\loigiai{
Phân tích dữ kiện, chọn định luật phù hợp. $M=Fd=17\cdot0,3=5,1$ N.m.
}
\end{bt}

% ===== Câu 51 =====
% ID: AIQ_L10C3_FULL_1347
% Mức: VD
\begin{bt}
% ID: AIQ_L10C3_FULL_1347
% Mức: VD
Lực 19 $N$ có giá cách trục quay $0,35\,\text{m}$. Tính mômen lực.
\loigiai{
Phân tích dữ kiện, chọn định luật phù hợp. $M=Fd=19\cdot0,35=6,65$ N.m.
}
\end{bt}

% ===== Câu 52 =====
% ID: AIQ_L10C3_FULL_1348
% Mức: VD
\begin{bt}
% ID: AIQ_L10C3_FULL_1348
% Mức: VD
Lực 13 $N$ có giá cách trục quay $0,4\,\text{m}$. Tính mômen lực.
\loigiai{
Phân tích dữ kiện, chọn định luật phù hợp. $M=Fd=13\cdot0,4=5,2$ N.m.
}
\end{bt}

% ===== Câu 53 =====
% ID: AIQ_L10C3_FULL_1349
% Mức: VD
\begin{bt}
% ID: AIQ_L10C3_FULL_1349
% Mức: VD
Lực 15 $N$ có giá cách trục quay $0,45\,\text{m}$. Tính mômen lực.
\loigiai{
Phân tích dữ kiện, chọn định luật phù hợp. $M=Fd=15\cdot0,45=6,75$ N.m.
}
\end{bt}

% ===== Câu 54 =====
% ID: AIQ_L10C3_FULL_1350
% Mức: VDC
\begin{bt}
% ID: AIQ_L10C3_FULL_1350
% Mức: VDC
Lực 17 $N$ có giá cách trục quay $0,2\,\text{m}$. Tính mômen lực.
\loigiai{
Phân tích dữ kiện, chọn định luật phù hợp. $M=Fd=17\cdot0,2=3,4$ N.m.
}
\end{bt}

% ===== Câu 55 =====
% ID: AIQ_L10C3_FULL_1351
% Mức: NB
\dangbt{Điều kiện cân bằng của vật có trục quay cố định}
\begin{ex}
% ID: AIQ_L10C3_FULL_1351
% Mức: NB
Hai lực gây mômen ngược chiều. $F_1=19N$, $d_1=0,35$ m, $d_2=0,45$ m. Tính $F_2$ để cân bằng.
\choice
{16,78 $N$}
{\True 14,78 $N$}
{17,74 $N$}
{11,82 $N$}
\loigiai{
Cân bằng: $F_1d_1=F_2d_2$, nên $F_2=14,78$ N. Chọn B.
}
\end{ex}

% ===== Câu 56 =====
% ID: AIQ_L10C3_FULL_1352
% Mức: NB
\begin{ex}
% ID: AIQ_L10C3_FULL_1352
% Mức: NB
Hai lực gây mômen ngược chiều. $F_1=13N$, $d_1=0,2$ m, $d_2=0,5$ m. Tính $F_2$ để cân bằng.
\choice
{\True 5,2 $N$}
{6,24 $N$}
{4,16 $N$}
{7,2 $N$}
\loigiai{
Cân bằng: $F_1d_1=F_2d_2$, nên $F_2=5,2$ N. Chọn A.
}
\end{ex}

% ===== Câu 57 =====
% ID: AIQ_L10C3_FULL_1353
% Mức: NB
\begin{ex}
% ID: AIQ_L10C3_FULL_1353
% Mức: NB
Hai lực gây mômen ngược chiều. $F_1=15N$, $d_1=0,25$ m, $d_2=0,55$ m. Tính $F_2$ để cân bằng.
\choice
{8,18 $N$}
{5,46 $N$}
{8,82 $N$}
{\True 6,82 $N$}
\loigiai{
Cân bằng: $F_1d_1=F_2d_2$, nên $F_2=6,82$ N. Chọn D.
}
\end{ex}

% ===== Câu 58 =====
% ID: AIQ_L10C3_FULL_1354
% Mức: TH
\begin{ex}
% ID: AIQ_L10C3_FULL_1354
% Mức: TH
Hai lực gây mômen ngược chiều. $F_1=17N$, $d_1=0,3$ m, $d_2=0,6$ m. Tính $F_2$ để cân bằng.
\choice
{6,8 $N$}
{10,5 $N$}
{\True 8,5 $N$}
{10,2 $N$}
\loigiai{
Cân bằng: $F_1d_1=F_2d_2$, nên $F_2=8,5$ N. Chọn C.
}
\end{ex}

% ===== Câu 59 =====
% ID: AIQ_L10C3_FULL_1355
% Mức: TH
\begin{ex}
% ID: AIQ_L10C3_FULL_1355
% Mức: TH
Hai lực gây mômen ngược chiều. $F_1=19N$, $d_1=0,35$ m, $d_2=0,4$ m. Tính $F_2$ để cân bằng.
\choice
{18,62 $N$}
{\True 16,62 $N$}
{19,94 $N$}
{13,3 $N$}
\loigiai{
Cân bằng: $F_1d_1=F_2d_2$, nên $F_2=16,62$ N. Chọn B.
}
\end{ex}

% ===== Câu 60 =====
% ID: AIQ_L10C3_FULL_1356
% Mức: TH
\begin{ex}
% ID: AIQ_L10C3_FULL_1356
% Mức: TH
Hai lực gây mômen ngược chiều. $F_1=13N$, $d_1=0,2$ m, $d_2=0,45$ m. Tính $F_2$ để cân bằng.
\choice
{\True 5,78 $N$}
{6,94 $N$}
{4,62 $N$}
{7,78 $N$}
\loigiai{
Cân bằng: $F_1d_1=F_2d_2$, nên $F_2=5,78$ N. Chọn A.
}
\end{ex}

% ===== Câu 61 =====
% ID: AIQ_L10C3_FULL_1357
% Mức: TH
\begin{ex}
% ID: AIQ_L10C3_FULL_1357
% Mức: TH
Hai lực gây mômen ngược chiều. $F_1=15N$, $d_1=0,25$ m, $d_2=0,5$ m. Tính $F_2$ để cân bằng.
\choice
{9 $N$}
{6 $N$}
{9,5 $N$}
{\True 7,5 $N$}
\loigiai{
Cân bằng: $F_1d_1=F_2d_2$, nên $F_2=7,5$ N. Chọn D.
}
\end{ex}

% ===== Câu 62 =====
% ID: AIQ_L10C3_FULL_1358
% Mức: VD
\begin{ex}
% ID: AIQ_L10C3_FULL_1358
% Mức: VD
Hai lực gây mômen ngược chiều. $F_1=17N$, $d_1=0,3$ m, $d_2=0,55$ m. Tính $F_2$ để cân bằng.
\choice
{7,42 $N$}
{11,27 $N$}
{\True 9,27 $N$}
{11,12 $N$}
\loigiai{
Cân bằng: $F_1d_1=F_2d_2$, nên $F_2=9,27$ N. Chọn C.
}
\end{ex}

% ===== Câu 63 =====
% ID: AIQ_L10C3_FULL_1359
% Mức: VD
\begin{ex}
% ID: AIQ_L10C3_FULL_1359
% Mức: VD
Hai lực gây mômen ngược chiều. $F_1=19N$, $d_1=0,35$ m, $d_2=0,6$ m. Tính $F_2$ để cân bằng.
\choice
{13,08 $N$}
{\True 11,08 $N$}
{13,3 $N$}
{8,86 $N$}
\loigiai{
Cân bằng: $F_1d_1=F_2d_2$, nên $F_2=11,08$ N. Chọn B.
}
\end{ex}

% ===== Câu 64 =====
% ID: AIQ_L10C3_FULL_1360
% Mức: VDC
\begin{ex}
% ID: AIQ_L10C3_FULL_1360
% Mức: VDC
Hai lực gây mômen ngược chiều. $F_1=13N$, $d_1=0,2$ m, $d_2=0,4$ m. Tính $F_2$ để cân bằng.
\choice
{\True 6,5 $N$}
{7,8 $N$}
{5,2 $N$}
{8,5 $N$}
\loigiai{
Cân bằng: $F_1d_1=F_2d_2$, nên $F_2=6,5$ N. Chọn A.
}
\end{ex}

% ===== Câu 65 =====
% ID: AIQ_L10C3_FULL_1361
% Mức: NB
\begin{ex}
% ID: AIQ_L10C3_FULL_1361
% Mức: NB
Hai lực gây mômen ngược chiều. $F_1=15N$, $d_1=0,25$ m, $d_2=0,45$ m. Tính $F_2$ để cân bằng.
\choiceTF
{\True Kết quả đúng là $8,33\,\text{N}$.}
{Kết quả bằng $16,66\,\text{N}$.}
{\True Khi hợp lực bằng 0, vật có gia tốc bằng 0.}
{Có thể bỏ qua đơn vị của đại lượng trong kết quả vật lí.}
\loigiai{
\begin{itemchoice}
\item (Đúng) Kết quả đúng là $8,33\,\text{N}$. (Vì): Cân bằng: $F_1d_1=F_2d_2$, nên $F_2=8,33$ N. b) Sai vì giá trị hoặc chiều nêu ra không phù hợp. c) Đúng theo định luật II Newton. d) Sai vì kết quả vật lí phải kèm đơn vị thích hợp
\item (Sai) Kết quả bằng $16,66\,\text{N}$.
\item (Đúng) Khi hợp lực bằng 0, vật có gia tốc bằng 0.
\item (Sai) Có thể bỏ qua đơn vị của đại lượng trong kết quả vật lí.
\end{itemchoice}
}
\end{ex}

% ===== Câu 66 =====
% ID: AIQ_L10C3_FULL_1362
% Mức: NB
\begin{ex}
% ID: AIQ_L10C3_FULL_1362
% Mức: NB
Hai lực gây mômen ngược chiều. $F_1=17N$, $d_1=0,3$ m, $d_2=0,5$ m. Tính $F_2$ để cân bằng.
\choiceTF
{\True Kết quả đúng là $10,2\,\text{N}$.}
{Kết quả bằng $20,4\,\text{N}$.}
{\True Khi hợp lực bằng 0, vật có gia tốc bằng 0.}
{Có thể bỏ qua đơn vị của đại lượng trong kết quả vật lí.}
\loigiai{
\begin{itemchoice}
\item (Đúng) Kết quả đúng là $10,2\,\text{N}$. (Vì): Cân bằng: $F_1d_1=F_2d_2$, nên $F_2=10,2$ N. b) Sai vì giá trị hoặc chiều nêu ra không phù hợp. c) Đúng theo định luật II Newton. d) Sai vì kết quả vật lí phải kèm đơn vị thích hợp
\item (Sai) Kết quả bằng $20,4\,\text{N}$.
\item (Đúng) Khi hợp lực bằng 0, vật có gia tốc bằng 0.
\item (Sai) Có thể bỏ qua đơn vị của đại lượng trong kết quả vật lí.
\end{itemchoice}
}
\end{ex}

% ===== Câu 67 =====
% ID: AIQ_L10C3_FULL_1363
% Mức: TH
\begin{ex}
% ID: AIQ_L10C3_FULL_1363
% Mức: TH
Hai lực gây mômen ngược chiều. $F_1=19N$, $d_1=0,35$ m, $d_2=0,55$ m. Tính $F_2$ để cân bằng.
\choiceTF
{\True Kết quả đúng là $12,09\,\text{N}$.}
{Kết quả bằng $24,18\,\text{N}$.}
{\True Khi hợp lực bằng 0, vật có gia tốc bằng 0.}
{Có thể bỏ qua đơn vị của đại lượng trong kết quả vật lí.}
\loigiai{
\begin{itemchoice}
\item (Đúng) Kết quả đúng là $12,09\,\text{N}$. (Vì): Cân bằng: $F_1d_1=F_2d_2$, nên $F_2=12,09$ N. b) Sai vì giá trị hoặc chiều nêu ra không phù hợp. c) Đúng theo định luật II Newton. d) Sai vì kết quả vật lí phải kèm đơn vị thích hợp
\item (Sai) Kết quả bằng $24,18\,\text{N}$.
\item (Đúng) Khi hợp lực bằng 0, vật có gia tốc bằng 0.
\item (Sai) Có thể bỏ qua đơn vị của đại lượng trong kết quả vật lí.
\end{itemchoice}
}
\end{ex}

% ===== Câu 68 =====
% ID: AIQ_L10C3_FULL_1364
% Mức: TH
\begin{ex}
% ID: AIQ_L10C3_FULL_1364
% Mức: TH
Hai lực gây mômen ngược chiều. $F_1=13N$, $d_1=0,2$ m, $d_2=0,6$ m. Tính $F_2$ để cân bằng.
\choiceTF
{\True Kết quả đúng là $4,33\,\text{N}$.}
{Kết quả bằng $8,66\,\text{N}$.}
{\True Khi hợp lực bằng 0, vật có gia tốc bằng 0.}
{Có thể bỏ qua đơn vị của đại lượng trong kết quả vật lí.}
\loigiai{
\begin{itemchoice}
\item (Đúng) Kết quả đúng là $4,33\,\text{N}$. (Vì): Cân bằng: $F_1d_1=F_2d_2$, nên $F_2=4,33$ N. b) Sai vì giá trị hoặc chiều nêu ra không phù hợp. c) Đúng theo định luật II Newton. d) Sai vì kết quả vật lí phải kèm đơn vị thích hợp
\item (Sai) Kết quả bằng $8,66\,\text{N}$.
\item (Đúng) Khi hợp lực bằng 0, vật có gia tốc bằng 0.
\item (Sai) Có thể bỏ qua đơn vị của đại lượng trong kết quả vật lí.
\end{itemchoice}
}
\end{ex}

% ===== Câu 69 =====
% ID: AIQ_L10C3_FULL_1365
% Mức: VD
\begin{ex}
% ID: AIQ_L10C3_FULL_1365
% Mức: VD
Hai lực gây mômen ngược chiều. $F_1=15N$, $d_1=0,25$ m, $d_2=0,4$ m. Tính $F_2$ để cân bằng.
\choiceTF
{\True Kết quả đúng là $9,38\,\text{N}$.}
{Kết quả bằng $18,76\,\text{N}$.}
{\True Khi hợp lực bằng 0, vật có gia tốc bằng 0.}
{Có thể bỏ qua đơn vị của đại lượng trong kết quả vật lí.}
\loigiai{
\begin{itemchoice}
\item (Đúng) Kết quả đúng là $9,38\,\text{N}$. (Vì): Cân bằng: $F_1d_1=F_2d_2$, nên $F_2=9,38$ N. b) Sai vì giá trị hoặc chiều nêu ra không phù hợp. c) Đúng theo định luật II Newton. d) Sai vì kết quả vật lí phải kèm đơn vị thích hợp
\item (Sai) Kết quả bằng $18,76\,\text{N}$.
\item (Đúng) Khi hợp lực bằng 0, vật có gia tốc bằng 0.
\item (Sai) Có thể bỏ qua đơn vị của đại lượng trong kết quả vật lí.
\end{itemchoice}
}
\end{ex}

% ===== Câu 70 =====
% ID: AIQ_L10C3_FULL_1366
% Mức: TH
\begin{ex}
% ID: AIQ_L10C3_FULL_1366
% Mức: TH
Hai lực gây mômen ngược chiều. $F_1=17N$, $d_1=0,3$ m, $d_2=0,45$ m. Tính $F_2$ để cân bằng.
\shortans{11,33}
\loigiai{
Cân bằng: $F_1d_1=F_2d_2$, nên $F_2=11,33$ N. Đáp án trả lời ngắn không ghi đơn vị.
}
\end{ex}

% ===== Câu 71 =====
% ID: AIQ_L10C3_FULL_1367
% Mức: TH
\begin{ex}
% ID: AIQ_L10C3_FULL_1367
% Mức: TH
Hai lực gây mômen ngược chiều. $F_1=19N$, $d_1=0,35$ m, $d_2=0,5$ m. Tính $F_2$ để cân bằng.
\shortans{13,3}
\loigiai{
Cân bằng: $F_1d_1=F_2d_2$, nên $F_2=13,3$ N. Đáp án trả lời ngắn không ghi đơn vị.
}
\end{ex}

% ===== Câu 72 =====
% ID: AIQ_L10C3_FULL_1368
% Mức: VD
\begin{ex}
% ID: AIQ_L10C3_FULL_1368
% Mức: VD
Hai lực gây mômen ngược chiều. $F_1=13N$, $d_1=0,2$ m, $d_2=0,55$ m. Tính $F_2$ để cân bằng.
\shortans{4,73}
\loigiai{
Cân bằng: $F_1d_1=F_2d_2$, nên $F_2=4,73$ N. Đáp án trả lời ngắn không ghi đơn vị.
}
\end{ex}

% ===== Câu 73 =====
% ID: AIQ_L10C3_FULL_1369
% Mức: VD
\begin{ex}
% ID: AIQ_L10C3_FULL_1369
% Mức: VD
Hai lực gây mômen ngược chiều. $F_1=15N$, $d_1=0,25$ m, $d_2=0,6$ m. Tính $F_2$ để cân bằng.
\shortans{6,25}
\loigiai{
Cân bằng: $F_1d_1=F_2d_2$, nên $F_2=6,25$ N. Đáp án trả lời ngắn không ghi đơn vị.
}
\end{ex}

% ===== Câu 74 =====
% ID: AIQ_L10C3_FULL_1370
% Mức: VD
\begin{ex}
% ID: AIQ_L10C3_FULL_1370
% Mức: VD
Hai lực gây mômen ngược chiều. $F_1=17N$, $d_1=0,3$ m, $d_2=0,4$ m. Tính $F_2$ để cân bằng.
\shortans{12,75}
\loigiai{
Cân bằng: $F_1d_1=F_2d_2$, nên $F_2=12,75$ N. Đáp án trả lời ngắn không ghi đơn vị.
}
\end{ex}

% ===== Câu 75 =====
% ID: AIQ_L10C3_FULL_1371
% Mức: VDC
\begin{ex}
% ID: AIQ_L10C3_FULL_1371
% Mức: VDC
Hai lực gây mômen ngược chiều. $F_1=19N$, $d_1=0,35$ m, $d_2=0,45$ m. Tính $F_2$ để cân bằng.
\shortans{14,78}
\loigiai{
Cân bằng: $F_1d_1=F_2d_2$, nên $F_2=14,78$ N. Đáp án trả lời ngắn không ghi đơn vị.
}
\end{ex}

% ===== Câu 76 =====
% ID: AIQ_L10C3_FULL_1372
% Mức: TH
\begin{bt}
% ID: AIQ_L10C3_FULL_1372
% Mức: TH
Hai lực gây mômen ngược chiều. $F_1=13N$, $d_1=0,2$ m, $d_2=0,5$ m. Tính $F_2$ để cân bằng.
\loigiai{
Phân tích dữ kiện, chọn định luật phù hợp. Cân bằng: $F_1d_1=F_2d_2$, nên $F_2=5,2$ N.
}
\end{bt}

% ===== Câu 77 =====
% ID: AIQ_L10C3_FULL_1373
% Mức: TH
\begin{bt}
% ID: AIQ_L10C3_FULL_1373
% Mức: TH
Hai lực gây mômen ngược chiều. $F_1=15N$, $d_1=0,25$ m, $d_2=0,55$ m. Tính $F_2$ để cân bằng.
\loigiai{
Phân tích dữ kiện, chọn định luật phù hợp. Cân bằng: $F_1d_1=F_2d_2$, nên $F_2=6,82$ N.
}
\end{bt}

% ===== Câu 78 =====
% ID: AIQ_L10C3_FULL_1374
% Mức: VD
\begin{bt}
% ID: AIQ_L10C3_FULL_1374
% Mức: VD
Hai lực gây mômen ngược chiều. $F_1=17N$, $d_1=0,3$ m, $d_2=0,6$ m. Tính $F_2$ để cân bằng.
\loigiai{
Phân tích dữ kiện, chọn định luật phù hợp. Cân bằng: $F_1d_1=F_2d_2$, nên $F_2=8,5$ N.
}
\end{bt}

% ===== Câu 79 =====
% ID: AIQ_L10C3_FULL_1375
% Mức: VD
\begin{bt}
% ID: AIQ_L10C3_FULL_1375
% Mức: VD
Hai lực gây mômen ngược chiều. $F_1=19N$, $d_1=0,35$ m, $d_2=0,4$ m. Tính $F_2$ để cân bằng.
\loigiai{
Phân tích dữ kiện, chọn định luật phù hợp. Cân bằng: $F_1d_1=F_2d_2$, nên $F_2=16,62$ N.
}
\end{bt}

% ===== Câu 80 =====
% ID: AIQ_L10C3_FULL_1376
% Mức: VD
\begin{bt}
% ID: AIQ_L10C3_FULL_1376
% Mức: VD
Hai lực gây mômen ngược chiều. $F_1=13N$, $d_1=0,2$ m, $d_2=0,45$ m. Tính $F_2$ để cân bằng.
\loigiai{
Phân tích dữ kiện, chọn định luật phù hợp. Cân bằng: $F_1d_1=F_2d_2$, nên $F_2=5,78$ N.
}
\end{bt}

% ===== Câu 81 =====
% ID: AIQ_L10C3_FULL_1377
% Mức: VDC
\begin{bt}
% ID: AIQ_L10C3_FULL_1377
% Mức: VDC
Hai lực gây mômen ngược chiều. $F_1=15N$, $d_1=0,25$ m, $d_2=0,5$ m. Tính $F_2$ để cân bằng.
\loigiai{
Phân tích dữ kiện, chọn định luật phù hợp. Cân bằng: $F_1d_1=F_2d_2$, nên $F_2=7,5$ N.
}
\end{bt}

% ===== Câu 82 =====
% ID: AIQ_L10C3_FULL_1378
% Mức: NB
\dangbt{Quy tắc mômen lực và bài toán đòn bẩy}
\begin{ex}
% ID: AIQ_L10C3_FULL_1378
% Mức: NB
Dùng đòn bẩy nâng vật nặng $140\,\text{N}$. Cánh tay đòn của vật là $0,16\,\text{m}$, của lực tác dụng là $0,7\,\text{m}$. Bỏ qua trọng lượng đòn. Tính lực cần tác dụng.
\choice
{25,6 $N$}
{34 $N$}
{\True 32 $N$}
{38,4 $N$}
\loigiai{
Theo quy tắc mômen: $F\,0,7=140\,0,16$, suy ra $F=32$ N. Chọn C.
}
\end{ex}

% ===== Câu 83 =====
% ID: AIQ_L10C3_FULL_1379
% Mức: NB
\begin{ex}
% ID: AIQ_L10C3_FULL_1379
% Mức: NB
Dùng đòn bẩy nâng vật nặng $160\,\text{N}$. Cánh tay đòn của vật là $0,18\,\text{m}$, của lực tác dụng là $0,75\,\text{m}$. Bỏ qua trọng lượng đòn. Tính lực cần tác dụng.
\choice
{40,4 $N$}
{\True 38,4 $N$}
{46,08 $N$}
{30,72 $N$}
\loigiai{
Theo quy tắc mômen: $F\,0,75=160\,0,18$, suy ra $F=38,4$ N. Chọn B.
}
\end{ex}

% ===== Câu 84 =====
% ID: AIQ_L10C3_FULL_1380
% Mức: NB
\begin{ex}
% ID: AIQ_L10C3_FULL_1380
% Mức: NB
Dùng đòn bẩy nâng vật nặng $180\,\text{N}$. Cánh tay đòn của vật là $0,1\,\text{m}$, của lực tác dụng là $0,5\,\text{m}$. Bỏ qua trọng lượng đòn. Tính lực cần tác dụng.
\choice
{\True 36 $N$}
{43,2 $N$}
{28,8 $N$}
{38 $N$}
\loigiai{
Theo quy tắc mômen: $F\,0,5=180\,0,1$, suy ra $F=36$ N. Chọn A.
}
\end{ex}

% ===== Câu 85 =====
% ID: AIQ_L10C3_FULL_1381
% Mức: TH
\begin{ex}
% ID: AIQ_L10C3_FULL_1381
% Mức: TH
Dùng đòn bẩy nâng vật nặng $200\,\text{N}$. Cánh tay đòn của vật là $0,12\,\text{m}$, của lực tác dụng là $0,55\,\text{m}$. Bỏ qua trọng lượng đòn. Tính lực cần tác dụng.
\choice
{52,37 $N$}
{34,91 $N$}
{45,64 $N$}
{\True 43,64 $N$}
\loigiai{
Theo quy tắc mômen: $F\,0,55=200\,0,12$, suy ra $F=43,64$ N. Chọn D.
}
\end{ex}

% ===== Câu 86 =====
% ID: AIQ_L10C3_FULL_1382
% Mức: TH
\begin{ex}
% ID: AIQ_L10C3_FULL_1382
% Mức: TH
Dùng đòn bẩy nâng vật nặng $220\,\text{N}$. Cánh tay đòn của vật là $0,14\,\text{m}$, của lực tác dụng là $0,6\,\text{m}$. Bỏ qua trọng lượng đòn. Tính lực cần tác dụng.
\choice
{41,06 $N$}
{53,33 $N$}
{\True 51,33 $N$}
{61,6 $N$}
\loigiai{
Theo quy tắc mômen: $F\,0,6=220\,0,14$, suy ra $F=51,33$ N. Chọn C.
}
\end{ex}

% ===== Câu 87 =====
% ID: AIQ_L10C3_FULL_1383
% Mức: TH
\begin{ex}
% ID: AIQ_L10C3_FULL_1383
% Mức: TH
Dùng đòn bẩy nâng vật nặng $240\,\text{N}$. Cánh tay đòn của vật là $0,16\,\text{m}$, của lực tác dụng là $0,65\,\text{m}$. Bỏ qua trọng lượng đòn. Tính lực cần tác dụng.
\choice
{61,08 $N$}
{\True 59,08 $N$}
{70,9 $N$}
{47,26 $N$}
\loigiai{
Theo quy tắc mômen: $F\,0,65=240\,0,16$, suy ra $F=59,08$ N. Chọn B.
}
\end{ex}

% ===== Câu 88 =====
% ID: AIQ_L10C3_FULL_1384
% Mức: TH
\begin{ex}
% ID: AIQ_L10C3_FULL_1384
% Mức: TH
Dùng đòn bẩy nâng vật nặng $100\,\text{N}$. Cánh tay đòn của vật là $0,18\,\text{m}$, của lực tác dụng là $0,7\,\text{m}$. Bỏ qua trọng lượng đòn. Tính lực cần tác dụng.
\choice
{\True 25,71 $N$}
{30,85 $N$}
{20,57 $N$}
{27,71 $N$}
\loigiai{
Theo quy tắc mômen: $F\,0,7=100\,0,18$, suy ra $F=25,71$ N. Chọn A.
}
\end{ex}

% ===== Câu 89 =====
% ID: AIQ_L10C3_FULL_1385
% Mức: VD
\begin{ex}
% ID: AIQ_L10C3_FULL_1385
% Mức: VD
Dùng đòn bẩy nâng vật nặng $120\,\text{N}$. Cánh tay đòn của vật là $0,1\,\text{m}$, của lực tác dụng là $0,75\,\text{m}$. Bỏ qua trọng lượng đòn. Tính lực cần tác dụng.
\choice
{19,2 $N$}
{12,8 $N$}
{18 $N$}
{\True 16 $N$}
\loigiai{
Theo quy tắc mômen: $F\,0,75=120\,0,1$, suy ra $F=16$ N. Chọn D.
}
\end{ex}

% ===== Câu 90 =====
% ID: AIQ_L10C3_FULL_1386
% Mức: VD
\begin{ex}
% ID: AIQ_L10C3_FULL_1386
% Mức: VD
Dùng đòn bẩy nâng vật nặng $140\,\text{N}$. Cánh tay đòn của vật là $0,12\,\text{m}$, của lực tác dụng là $0,5\,\text{m}$. Bỏ qua trọng lượng đòn. Tính lực cần tác dụng.
\choice
{26,88 $N$}
{35,6 $N$}
{\True 33,6 $N$}
{40,32 $N$}
\loigiai{
Theo quy tắc mômen: $F\,0,5=140\,0,12$, suy ra $F=33,6$ N. Chọn C.
}
\end{ex}

% ===== Câu 91 =====
% ID: AIQ_L10C3_FULL_1387
% Mức: VDC
\begin{ex}
% ID: AIQ_L10C3_FULL_1387
% Mức: VDC
Dùng đòn bẩy nâng vật nặng $160\,\text{N}$. Cánh tay đòn của vật là $0,14\,\text{m}$, của lực tác dụng là $0,55\,\text{m}$. Bỏ qua trọng lượng đòn. Tính lực cần tác dụng.
\choice
{42,73 $N$}
{\True 40,73 $N$}
{48,88 $N$}
{32,58 $N$}
\loigiai{
Theo quy tắc mômen: $F\,0,55=160\,0,14$, suy ra $F=40,73$ N. Chọn B.
}
\end{ex}

% ===== Câu 92 =====
% ID: AIQ_L10C3_FULL_1388
% Mức: NB
\begin{ex}
% ID: AIQ_L10C3_FULL_1388
% Mức: NB
Dùng đòn bẩy nâng vật nặng $180\,\text{N}$. Cánh tay đòn của vật là $0,16\,\text{m}$, của lực tác dụng là $0,6\,\text{m}$. Bỏ qua trọng lượng đòn. Tính lực cần tác dụng.
\choiceTF
{\True Kết quả đúng là $48\,\text{N}$.}
{Kết quả bằng $96\,\text{N}$.}
{\True Khi hợp lực bằng 0, vật có gia tốc bằng 0.}
{Có thể bỏ qua đơn vị của đại lượng trong kết quả vật lí.}
\loigiai{
\begin{itemchoice}
\item (Đúng) Kết quả đúng là $48\,\text{N}$. (Vì): Theo quy tắc mômen: $F\,0,6=180\,0,16$, suy ra $F=48$ N. b) Sai vì giá trị hoặc chiều nêu ra không phù hợp. c) Đúng theo định luật II Newton. d) Sai vì kết quả vật lí phải kèm đơn vị thích hợp
\item (Sai) Kết quả bằng $96\,\text{N}$.
\item (Đúng) Khi hợp lực bằng 0, vật có gia tốc bằng 0.
\item (Sai) Có thể bỏ qua đơn vị của đại lượng trong kết quả vật lí.
\end{itemchoice}
}
\end{ex}

% ===== Câu 93 =====
% ID: AIQ_L10C3_FULL_1389
% Mức: NB
\begin{ex}
% ID: AIQ_L10C3_FULL_1389
% Mức: NB
Dùng đòn bẩy nâng vật nặng $200\,\text{N}$. Cánh tay đòn của vật là $0,18\,\text{m}$, của lực tác dụng là $0,65\,\text{m}$. Bỏ qua trọng lượng đòn. Tính lực cần tác dụng.
\choiceTF
{\True Kết quả đúng là $55,38\,\text{N}$.}
{Kết quả bằng $110,76\,\text{N}$.}
{\True Khi hợp lực bằng 0, vật có gia tốc bằng 0.}
{Có thể bỏ qua đơn vị của đại lượng trong kết quả vật lí.}
\loigiai{
\begin{itemchoice}
\item (Đúng) Kết quả đúng là $55,38\,\text{N}$. (Vì): Theo quy tắc mômen: $F\,0,65=200\,0,18$, suy ra $F=55,38$ N. b) Sai vì giá trị hoặc chiều nêu ra không phù hợp. c) Đúng theo định luật II Newton. d) Sai vì kết quả vật lí phải kèm đơn vị thích hợp
\item (Sai) Kết quả bằng $110,76\,\text{N}$.
\item (Đúng) Khi hợp lực bằng 0, vật có gia tốc bằng 0.
\item (Sai) Có thể bỏ qua đơn vị của đại lượng trong kết quả vật lí.
\end{itemchoice}
}
\end{ex}

% ===== Câu 94 =====
% ID: AIQ_L10C3_FULL_1390
% Mức: TH
\begin{ex}
% ID: AIQ_L10C3_FULL_1390
% Mức: TH
Dùng đòn bẩy nâng vật nặng $220\,\text{N}$. Cánh tay đòn của vật là $0,1\,\text{m}$, của lực tác dụng là $0,7\,\text{m}$. Bỏ qua trọng lượng đòn. Tính lực cần tác dụng.
\choiceTF
{\True Kết quả đúng là $31,43\,\text{N}$.}
{Kết quả bằng $62,86\,\text{N}$.}
{\True Khi hợp lực bằng 0, vật có gia tốc bằng 0.}
{Có thể bỏ qua đơn vị của đại lượng trong kết quả vật lí.}
\loigiai{
\begin{itemchoice}
\item (Đúng) Kết quả đúng là $31,43\,\text{N}$. (Vì): Theo quy tắc mômen: $F\,0,7=220\,0,1$, suy ra $F=31,43$ N. b) Sai vì giá trị hoặc chiều nêu ra không phù hợp. c) Đúng theo định luật II Newton. d) Sai vì kết quả vật lí phải kèm đơn vị thích hợp
\item (Sai) Kết quả bằng $62,86\,\text{N}$.
\item (Đúng) Khi hợp lực bằng 0, vật có gia tốc bằng 0.
\item (Sai) Có thể bỏ qua đơn vị của đại lượng trong kết quả vật lí.
\end{itemchoice}
}
\end{ex}

% ===== Câu 95 =====
% ID: AIQ_L10C3_FULL_1391
% Mức: TH
\begin{ex}
% ID: AIQ_L10C3_FULL_1391
% Mức: TH
Dùng đòn bẩy nâng vật nặng $240\,\text{N}$. Cánh tay đòn của vật là $0,12\,\text{m}$, của lực tác dụng là $0,75\,\text{m}$. Bỏ qua trọng lượng đòn. Tính lực cần tác dụng.
\choiceTF
{\True Kết quả đúng là $38,4\,\text{N}$.}
{Kết quả bằng $76,8\,\text{N}$.}
{\True Khi hợp lực bằng 0, vật có gia tốc bằng 0.}
{Có thể bỏ qua đơn vị của đại lượng trong kết quả vật lí.}
\loigiai{
\begin{itemchoice}
\item (Đúng) Kết quả đúng là $38,4\,\text{N}$. (Vì): Theo quy tắc mômen: $F\,0,75=240\,0,12$, suy ra $F=38,4$ N. b) Sai vì giá trị hoặc chiều nêu ra không phù hợp. c) Đúng theo định luật II Newton. d) Sai vì kết quả vật lí phải kèm đơn vị thích hợp
\item (Sai) Kết quả bằng $76,8\,\text{N}$.
\item (Đúng) Khi hợp lực bằng 0, vật có gia tốc bằng 0.
\item (Sai) Có thể bỏ qua đơn vị của đại lượng trong kết quả vật lí.
\end{itemchoice}
}
\end{ex}

% ===== Câu 96 =====
% ID: AIQ_L10C3_FULL_1392
% Mức: VD
\begin{ex}
% ID: AIQ_L10C3_FULL_1392
% Mức: VD
Dùng đòn bẩy nâng vật nặng $100\,\text{N}$. Cánh tay đòn của vật là $0,14\,\text{m}$, của lực tác dụng là $0,5\,\text{m}$. Bỏ qua trọng lượng đòn. Tính lực cần tác dụng.
\choiceTF
{\True Kết quả đúng là $28\,\text{N}$.}
{Kết quả bằng $56\,\text{N}$.}
{\True Khi hợp lực bằng 0, vật có gia tốc bằng 0.}
{Có thể bỏ qua đơn vị của đại lượng trong kết quả vật lí.}
\loigiai{
\begin{itemchoice}
\item (Đúng) Kết quả đúng là $28\,\text{N}$. (Vì): Theo quy tắc mômen: $F\,0,5=100\,0,14$, suy ra $F=28$ N. b) Sai vì giá trị hoặc chiều nêu ra không phù hợp. c) Đúng theo định luật II Newton. d) Sai vì kết quả vật lí phải kèm đơn vị thích hợp
\item (Sai) Kết quả bằng $56\,\text{N}$.
\item (Đúng) Khi hợp lực bằng 0, vật có gia tốc bằng 0.
\item (Sai) Có thể bỏ qua đơn vị của đại lượng trong kết quả vật lí.
\end{itemchoice}
}
\end{ex}

% ===== Câu 97 =====
% ID: AIQ_L10C3_FULL_1393
% Mức: TH
\begin{ex}
% ID: AIQ_L10C3_FULL_1393
% Mức: TH
Dùng đòn bẩy nâng vật nặng $120\,\text{N}$. Cánh tay đòn của vật là $0,16\,\text{m}$, của lực tác dụng là $0,55\,\text{m}$. Bỏ qua trọng lượng đòn. Tính lực cần tác dụng.
\shortans{34,91}
\loigiai{
Theo quy tắc mômen: $F\,0,55=120\,0,16$, suy ra $F=34,91$ N. Đáp án trả lời ngắn không ghi đơn vị.
}
\end{ex}

% ===== Câu 98 =====
% ID: AIQ_L10C3_FULL_1394
% Mức: TH
\begin{ex}
% ID: AIQ_L10C3_FULL_1394
% Mức: TH
Dùng đòn bẩy nâng vật nặng $140\,\text{N}$. Cánh tay đòn của vật là $0,18\,\text{m}$, của lực tác dụng là $0,6\,\text{m}$. Bỏ qua trọng lượng đòn. Tính lực cần tác dụng.
\shortans{42}
\loigiai{
Theo quy tắc mômen: $F\,0,6=140\,0,18$, suy ra $F=42$ N. Đáp án trả lời ngắn không ghi đơn vị.
}
\end{ex}

% ===== Câu 99 =====
% ID: AIQ_L10C3_FULL_1395
% Mức: VD
\begin{ex}
% ID: AIQ_L10C3_FULL_1395
% Mức: VD
Dùng đòn bẩy nâng vật nặng $160\,\text{N}$. Cánh tay đòn của vật là $0,1\,\text{m}$, của lực tác dụng là $0,65\,\text{m}$. Bỏ qua trọng lượng đòn. Tính lực cần tác dụng.
\shortans{24,62}
\loigiai{
Theo quy tắc mômen: $F\,0,65=160\,0,1$, suy ra $F=24,62$ N. Đáp án trả lời ngắn không ghi đơn vị.
}
\end{ex}

% ===== Câu 100 =====
% ID: AIQ_L10C3_FULL_1396
% Mức: VD
\begin{ex}
% ID: AIQ_L10C3_FULL_1396
% Mức: VD
Dùng đòn bẩy nâng vật nặng $180\,\text{N}$. Cánh tay đòn của vật là $0,12\,\text{m}$, của lực tác dụng là $0,7\,\text{m}$. Bỏ qua trọng lượng đòn. Tính lực cần tác dụng.
\shortans{30,86}
\loigiai{
Theo quy tắc mômen: $F\,0,7=180\,0,12$, suy ra $F=30,86$ N. Đáp án trả lời ngắn không ghi đơn vị.
}
\end{ex}

% ===== Câu 101 =====
% ID: AIQ_L10C3_FULL_1397
% Mức: VD
\begin{ex}
% ID: AIQ_L10C3_FULL_1397
% Mức: VD
Dùng đòn bẩy nâng vật nặng $200\,\text{N}$. Cánh tay đòn của vật là $0,14\,\text{m}$, của lực tác dụng là $0,75\,\text{m}$. Bỏ qua trọng lượng đòn. Tính lực cần tác dụng.
\shortans{37,33}
\loigiai{
Theo quy tắc mômen: $F\,0,75=200\,0,14$, suy ra $F=37,33$ N. Đáp án trả lời ngắn không ghi đơn vị.
}
\end{ex}

% ===== Câu 102 =====
% ID: AIQ_L10C3_FULL_1398
% Mức: VDC
\begin{ex}
% ID: AIQ_L10C3_FULL_1398
% Mức: VDC
Dùng đòn bẩy nâng vật nặng $220\,\text{N}$. Cánh tay đòn của vật là $0,16\,\text{m}$, của lực tác dụng là $0,5\,\text{m}$. Bỏ qua trọng lượng đòn. Tính lực cần tác dụng.
\shortans{70,4}
\loigiai{
Theo quy tắc mômen: $F\,0,5=220\,0,16$, suy ra $F=70,4$ N. Đáp án trả lời ngắn không ghi đơn vị.
}
\end{ex}

% ===== Câu 103 =====
% ID: AIQ_L10C3_FULL_1399
% Mức: TH
\begin{bt}
% ID: AIQ_L10C3_FULL_1399
% Mức: TH
Dùng đòn bẩy nâng vật nặng $240\,\text{N}$. Cánh tay đòn của vật là $0,18\,\text{m}$, của lực tác dụng là $0,55\,\text{m}$. Bỏ qua trọng lượng đòn. Tính lực cần tác dụng.
\loigiai{
Phân tích dữ kiện, chọn định luật phù hợp. Theo quy tắc mômen: $F\,0,55=240\,0,18$, suy ra $F=78,55$ N.
}
\end{bt}

% ===== Câu 104 =====
% ID: AIQ_L10C3_FULL_1400
% Mức: TH
\begin{bt}
% ID: AIQ_L10C3_FULL_1400
% Mức: TH
Dùng đòn bẩy nâng vật nặng $100\,\text{N}$. Cánh tay đòn của vật là $0,1\,\text{m}$, của lực tác dụng là $0,6\,\text{m}$. Bỏ qua trọng lượng đòn. Tính lực cần tác dụng.
\loigiai{
Phân tích dữ kiện, chọn định luật phù hợp. Theo quy tắc mômen: $F\,0,6=100\,0,1$, suy ra $F=16,67$ N.
}
\end{bt}

% ===== Câu 105 =====
% ID: AIQ_L10C3_FULL_1401
% Mức: VD
\begin{bt}
% ID: AIQ_L10C3_FULL_1401
% Mức: VD
Dùng đòn bẩy nâng vật nặng $120\,\text{N}$. Cánh tay đòn của vật là $0,12\,\text{m}$, của lực tác dụng là $0,65\,\text{m}$. Bỏ qua trọng lượng đòn. Tính lực cần tác dụng.
\loigiai{
Phân tích dữ kiện, chọn định luật phù hợp. Theo quy tắc mômen: $F\,0,65=120\,0,12$, suy ra $F=22,15$ N.
}
\end{bt}

% ===== Câu 106 =====
% ID: AIQ_L10C3_FULL_1402
% Mức: VD
\begin{bt}
% ID: AIQ_L10C3_FULL_1402
% Mức: VD
Dùng đòn bẩy nâng vật nặng $140\,\text{N}$. Cánh tay đòn của vật là $0,14\,\text{m}$, của lực tác dụng là $0,7\,\text{m}$. Bỏ qua trọng lượng đòn. Tính lực cần tác dụng.
\loigiai{
Phân tích dữ kiện, chọn định luật phù hợp. Theo quy tắc mômen: $F\,0,7=140\,0,14$, suy ra $F=28$ N.
}
\end{bt}

% ===== Câu 107 =====
% ID: AIQ_L10C3_FULL_1403
% Mức: VD
\begin{bt}
% ID: AIQ_L10C3_FULL_1403
% Mức: VD
Dùng đòn bẩy nâng vật nặng $160\,\text{N}$. Cánh tay đòn của vật là $0,16\,\text{m}$, của lực tác dụng là $0,75\,\text{m}$. Bỏ qua trọng lượng đòn. Tính lực cần tác dụng.
\loigiai{
Phân tích dữ kiện, chọn định luật phù hợp. Theo quy tắc mômen: $F\,0,75=160\,0,16$, suy ra $F=34,13$ N.
}
\end{bt}

% ===== Câu 108 =====
% ID: AIQ_L10C3_FULL_1404
% Mức: VDC
\begin{bt}
% ID: AIQ_L10C3_FULL_1404
% Mức: VDC
Dùng đòn bẩy nâng vật nặng $180\,\text{N}$. Cánh tay đòn của vật là $0,18\,\text{m}$, của lực tác dụng là $0,5\,\text{m}$. Bỏ qua trọng lượng đòn. Tính lực cần tác dụng.
\loigiai{
Phân tích dữ kiện, chọn định luật phù hợp. Theo quy tắc mômen: $F\,0,5=180\,0,18$, suy ra $F=64,8$ N.
}
\end{bt}

% ===== Câu 109 =====
% ID: AIQ_L10C3_FULL_1405
% Mức: NB
\dangbt{Cân bằng của thanh đồng chất chịu nhiều lực}
\begin{ex}
% ID: AIQ_L10C3_FULL_1405
% Mức: NB
Thanh đồng chất khối lượng $3\,\text{kg}$, dài $3\,\text{m}$ nằm ngang, được đỡ ở hai đầu. Tải phân bố đối xứng. Lấy g= $10\,\text{m}$ /s^2. Phản lực tại mỗi đầu bằng bao nhiêu?
\choice
{18 $N$}
{12 $N$}
{17 $N$}
{\True 15 $N$}
\loigiai{
Do đối xứng và cân bằng: $N_1=N_2=mg/2=15$ N. Chọn D.
}
\end{ex}

% ===== Câu 110 =====
% ID: AIQ_L10C3_FULL_1406
% Mức: NB
\begin{ex}
% ID: AIQ_L10C3_FULL_1406
% Mức: NB
Thanh đồng chất khối lượng $4\,\text{kg}$, dài $4\,\text{m}$ nằm ngang, được đỡ ở hai đầu. Tải phân bố đối xứng. Lấy g= $10\,\text{m}$ /s^2. Phản lực tại mỗi đầu bằng bao nhiêu?
\choice
{16 $N$}
{22 $N$}
{\True 20 $N$}
{24 $N$}
\loigiai{
Do đối xứng và cân bằng: $N_1=N_2=mg/2=20$ N. Chọn C.
}
\end{ex}

% ===== Câu 111 =====
% ID: AIQ_L10C3_FULL_1407
% Mức: NB
\begin{ex}
% ID: AIQ_L10C3_FULL_1407
% Mức: NB
Thanh đồng chất khối lượng $5\,\text{kg}$, dài $5\,\text{m}$ nằm ngang, được đỡ ở hai đầu. Tải phân bố đối xứng. Lấy g= $10\,\text{m}$ /s^2. Phản lực tại mỗi đầu bằng bao nhiêu?
\choice
{27 $N$}
{\True 25 $N$}
{30 $N$}
{20 $N$}
\loigiai{
Do đối xứng và cân bằng: $N_1=N_2=mg/2=25$ N. Chọn B.
}
\end{ex}

% ===== Câu 112 =====
% ID: AIQ_L10C3_FULL_1408
% Mức: TH
\begin{ex}
% ID: AIQ_L10C3_FULL_1408
% Mức: TH
Thanh đồng chất khối lượng $6\,\text{kg}$, dài $2\,\text{m}$ nằm ngang, được đỡ ở hai đầu. Tải phân bố đối xứng. Lấy g= $10\,\text{m}$ /s^2. Phản lực tại mỗi đầu bằng bao nhiêu?
\choice
{\True 30 $N$}
{36 $N$}
{24 $N$}
{32 $N$}
\loigiai{
Do đối xứng và cân bằng: $N_1=N_2=mg/2=30$ N. Chọn A.
}
\end{ex}

% ===== Câu 113 =====
% ID: AIQ_L10C3_FULL_1409
% Mức: TH
\begin{ex}
% ID: AIQ_L10C3_FULL_1409
% Mức: TH
Thanh đồng chất khối lượng $7\,\text{kg}$, dài $3\,\text{m}$ nằm ngang, được đỡ ở hai đầu. Tải phân bố đối xứng. Lấy g= $10\,\text{m}$ /s^2. Phản lực tại mỗi đầu bằng bao nhiêu?
\choice
{42 $N$}
{28 $N$}
{37 $N$}
{\True 35 $N$}
\loigiai{
Do đối xứng và cân bằng: $N_1=N_2=mg/2=35$ N. Chọn D.
}
\end{ex}

% ===== Câu 114 =====
% ID: AIQ_L10C3_FULL_1410
% Mức: TH
\begin{ex}
% ID: AIQ_L10C3_FULL_1410
% Mức: TH
Thanh đồng chất khối lượng $2\,\text{kg}$, dài $4\,\text{m}$ nằm ngang, được đỡ ở hai đầu. Tải phân bố đối xứng. Lấy g= $10\,\text{m}$ /s^2. Phản lực tại mỗi đầu bằng bao nhiêu?
\choice
{8 $N$}
{12 $N$}
{\True 10 $N$}
{12 $N$}
\loigiai{
Do đối xứng và cân bằng: $N_1=N_2=mg/2=10$ N. Chọn C.
}
\end{ex}

% ===== Câu 115 =====
% ID: AIQ_L10C3_FULL_1411
% Mức: TH
\begin{ex}
% ID: AIQ_L10C3_FULL_1411
% Mức: TH
Thanh đồng chất khối lượng $3\,\text{kg}$, dài $5\,\text{m}$ nằm ngang, được đỡ ở hai đầu. Tải phân bố đối xứng. Lấy g= $10\,\text{m}$ /s^2. Phản lực tại mỗi đầu bằng bao nhiêu?
\choice
{17 $N$}
{\True 15 $N$}
{18 $N$}
{12 $N$}
\loigiai{
Do đối xứng và cân bằng: $N_1=N_2=mg/2=15$ N. Chọn B.
}
\end{ex}

% ===== Câu 116 =====
% ID: AIQ_L10C3_FULL_1412
% Mức: VD
\begin{ex}
% ID: AIQ_L10C3_FULL_1412
% Mức: VD
Thanh đồng chất khối lượng $4\,\text{kg}$, dài $2\,\text{m}$ nằm ngang, được đỡ ở hai đầu. Tải phân bố đối xứng. Lấy g= $10\,\text{m}$ /s^2. Phản lực tại mỗi đầu bằng bao nhiêu?
\choice
{\True 20 $N$}
{24 $N$}
{16 $N$}
{22 $N$}
\loigiai{
Do đối xứng và cân bằng: $N_1=N_2=mg/2=20$ N. Chọn A.
}
\end{ex}

% ===== Câu 117 =====
% ID: AIQ_L10C3_FULL_1413
% Mức: VD
\begin{ex}
% ID: AIQ_L10C3_FULL_1413
% Mức: VD
Thanh đồng chất khối lượng $5\,\text{kg}$, dài $3\,\text{m}$ nằm ngang, được đỡ ở hai đầu. Tải phân bố đối xứng. Lấy g= $10\,\text{m}$ /s^2. Phản lực tại mỗi đầu bằng bao nhiêu?
\choice
{30 $N$}
{20 $N$}
{27 $N$}
{\True 25 $N$}
\loigiai{
Do đối xứng và cân bằng: $N_1=N_2=mg/2=25$ N. Chọn D.
}
\end{ex}

% ===== Câu 118 =====
% ID: AIQ_L10C3_FULL_1414
% Mức: VDC
\begin{ex}
% ID: AIQ_L10C3_FULL_1414
% Mức: VDC
Thanh đồng chất khối lượng $6\,\text{kg}$, dài $4\,\text{m}$ nằm ngang, được đỡ ở hai đầu. Tải phân bố đối xứng. Lấy g= $10\,\text{m}$ /s^2. Phản lực tại mỗi đầu bằng bao nhiêu?
\choice
{24 $N$}
{32 $N$}
{\True 30 $N$}
{36 $N$}
\loigiai{
Do đối xứng và cân bằng: $N_1=N_2=mg/2=30$ N. Chọn C.
}
\end{ex}

% ===== Câu 119 =====
% ID: AIQ_L10C3_FULL_1415
% Mức: NB
\begin{ex}
% ID: AIQ_L10C3_FULL_1415
% Mức: NB
Thanh đồng chất khối lượng $7\,\text{kg}$, dài $5\,\text{m}$ nằm ngang, được đỡ ở hai đầu. Tải phân bố đối xứng. Lấy g= $10\,\text{m}$ /s^2. Phản lực tại mỗi đầu bằng bao nhiêu?
\choiceTF
{\True Kết quả đúng là $35\,\text{N}$.}
{Kết quả bằng $70\,\text{N}$.}
{\True Khi hợp lực bằng 0, vật có gia tốc bằng 0.}
{Có thể bỏ qua đơn vị của đại lượng trong kết quả vật lí.}
\loigiai{
\begin{itemchoice}
\item (Đúng) Kết quả đúng là $35\,\text{N}$. (Vì): o đối xứng và cân bằng: $N_1=N_2=mg/2=35$ N. b) Sai vì giá trị hoặc chiều nêu ra không phù hợp. c) Đúng theo định luật II Newton. d) Sai vì kết quả vật lí phải kèm đơn vị thích hợp
\item (Sai) Kết quả bằng $70\,\text{N}$.
\item (Đúng) Khi hợp lực bằng 0, vật có gia tốc bằng 0.
\item (Sai) Có thể bỏ qua đơn vị của đại lượng trong kết quả vật lí.
\end{itemchoice}
}
\end{ex}

% ===== Câu 120 =====
% ID: AIQ_L10C3_FULL_1416
% Mức: NB
\begin{ex}
% ID: AIQ_L10C3_FULL_1416
% Mức: NB
Thanh đồng chất khối lượng $2\,\text{kg}$, dài $2\,\text{m}$ nằm ngang, được đỡ ở hai đầu. Tải phân bố đối xứng. Lấy g= $10\,\text{m}$ /s^2. Phản lực tại mỗi đầu bằng bao nhiêu?
\choiceTF
{\True Kết quả đúng là $10\,\text{N}$.}
{Kết quả bằng $20\,\text{N}$.}
{\True Khi hợp lực bằng 0, vật có gia tốc bằng 0.}
{Có thể bỏ qua đơn vị của đại lượng trong kết quả vật lí.}
\loigiai{
\begin{itemchoice}
\item (Đúng) Kết quả đúng là $10\,\text{N}$. (Vì): o đối xứng và cân bằng: $N_1=N_2=mg/2=10$ N. b) Sai vì giá trị hoặc chiều nêu ra không phù hợp. c) Đúng theo định luật II Newton. d) Sai vì kết quả vật lí phải kèm đơn vị thích hợp
\item (Sai) Kết quả bằng $20\,\text{N}$.
\item (Đúng) Khi hợp lực bằng 0, vật có gia tốc bằng 0.
\item (Sai) Có thể bỏ qua đơn vị của đại lượng trong kết quả vật lí.
\end{itemchoice}
}
\end{ex}

% ===== Câu 121 =====
% ID: AIQ_L10C3_FULL_1417
% Mức: TH
\begin{ex}
% ID: AIQ_L10C3_FULL_1417
% Mức: TH
Thanh đồng chất khối lượng $3\,\text{kg}$, dài $3\,\text{m}$ nằm ngang, được đỡ ở hai đầu. Tải phân bố đối xứng. Lấy g= $10\,\text{m}$ /s^2. Phản lực tại mỗi đầu bằng bao nhiêu?
\choiceTF
{\True Kết quả đúng là $15\,\text{N}$.}
{Kết quả bằng $30\,\text{N}$.}
{\True Khi hợp lực bằng 0, vật có gia tốc bằng 0.}
{Có thể bỏ qua đơn vị của đại lượng trong kết quả vật lí.}
\loigiai{
\begin{itemchoice}
\item (Đúng) Kết quả đúng là $15\,\text{N}$. (Vì): o đối xứng và cân bằng: $N_1=N_2=mg/2=15$ N. b) Sai vì giá trị hoặc chiều nêu ra không phù hợp. c) Đúng theo định luật II Newton. d) Sai vì kết quả vật lí phải kèm đơn vị thích hợp
\item (Sai) Kết quả bằng $30\,\text{N}$.
\item (Đúng) Khi hợp lực bằng 0, vật có gia tốc bằng 0.
\item (Sai) Có thể bỏ qua đơn vị của đại lượng trong kết quả vật lí.
\end{itemchoice}
}
\end{ex}

% ===== Câu 122 =====
% ID: AIQ_L10C3_FULL_1418
% Mức: TH
\begin{ex}
% ID: AIQ_L10C3_FULL_1418
% Mức: TH
Thanh đồng chất khối lượng $4\,\text{kg}$, dài $4\,\text{m}$ nằm ngang, được đỡ ở hai đầu. Tải phân bố đối xứng. Lấy g= $10\,\text{m}$ /s^2. Phản lực tại mỗi đầu bằng bao nhiêu?
\choiceTF
{\True Kết quả đúng là $20\,\text{N}$.}
{Kết quả bằng $40\,\text{N}$.}
{\True Khi hợp lực bằng 0, vật có gia tốc bằng 0.}
{Có thể bỏ qua đơn vị của đại lượng trong kết quả vật lí.}
\loigiai{
\begin{itemchoice}
\item (Đúng) Kết quả đúng là $20\,\text{N}$. (Vì): o đối xứng và cân bằng: $N_1=N_2=mg/2=20$ N. b) Sai vì giá trị hoặc chiều nêu ra không phù hợp. c) Đúng theo định luật II Newton. d) Sai vì kết quả vật lí phải kèm đơn vị thích hợp
\item (Sai) Kết quả bằng $40\,\text{N}$.
\item (Đúng) Khi hợp lực bằng 0, vật có gia tốc bằng 0.
\item (Sai) Có thể bỏ qua đơn vị của đại lượng trong kết quả vật lí.
\end{itemchoice}
}
\end{ex}

% ===== Câu 123 =====
% ID: AIQ_L10C3_FULL_1419
% Mức: VD
\begin{ex}
% ID: AIQ_L10C3_FULL_1419
% Mức: VD
Thanh đồng chất khối lượng $5\,\text{kg}$, dài $5\,\text{m}$ nằm ngang, được đỡ ở hai đầu. Tải phân bố đối xứng. Lấy g= $10\,\text{m}$ /s^2. Phản lực tại mỗi đầu bằng bao nhiêu?
\choiceTF
{\True Kết quả đúng là $25\,\text{N}$.}
{Kết quả bằng $50\,\text{N}$.}
{\True Khi hợp lực bằng 0, vật có gia tốc bằng 0.}
{Có thể bỏ qua đơn vị của đại lượng trong kết quả vật lí.}
\loigiai{
\begin{itemchoice}
\item (Đúng) Kết quả đúng là $25\,\text{N}$. (Vì): o đối xứng và cân bằng: $N_1=N_2=mg/2=25$ N. b) Sai vì giá trị hoặc chiều nêu ra không phù hợp. c) Đúng theo định luật II Newton. d) Sai vì kết quả vật lí phải kèm đơn vị thích hợp
\item (Sai) Kết quả bằng $50\,\text{N}$.
\item (Đúng) Khi hợp lực bằng 0, vật có gia tốc bằng 0.
\item (Sai) Có thể bỏ qua đơn vị của đại lượng trong kết quả vật lí.
\end{itemchoice}
}
\end{ex}

% ===== Câu 124 =====
% ID: AIQ_L10C3_FULL_1420
% Mức: TH
\begin{ex}
% ID: AIQ_L10C3_FULL_1420
% Mức: TH
Thanh đồng chất khối lượng $6\,\text{kg}$, dài $2\,\text{m}$ nằm ngang, được đỡ ở hai đầu. Tải phân bố đối xứng. Lấy g= $10\,\text{m}$ /s^2. Phản lực tại mỗi đầu bằng bao nhiêu?
\shortans{30}
\loigiai{
Do đối xứng và cân bằng: $N_1=N_2=mg/2=30$ N. Đáp án trả lời ngắn không ghi đơn vị.
}
\end{ex}

% ===== Câu 125 =====
% ID: AIQ_L10C3_FULL_1421
% Mức: TH
\begin{ex}
% ID: AIQ_L10C3_FULL_1421
% Mức: TH
Thanh đồng chất khối lượng $7\,\text{kg}$, dài $3\,\text{m}$ nằm ngang, được đỡ ở hai đầu. Tải phân bố đối xứng. Lấy g= $10\,\text{m}$ /s^2. Phản lực tại mỗi đầu bằng bao nhiêu?
\shortans{35}
\loigiai{
Do đối xứng và cân bằng: $N_1=N_2=mg/2=35$ N. Đáp án trả lời ngắn không ghi đơn vị.
}
\end{ex}

% ===== Câu 126 =====
% ID: AIQ_L10C3_FULL_1422
% Mức: VD
\begin{ex}
% ID: AIQ_L10C3_FULL_1422
% Mức: VD
Thanh đồng chất khối lượng $2\,\text{kg}$, dài $4\,\text{m}$ nằm ngang, được đỡ ở hai đầu. Tải phân bố đối xứng. Lấy g= $10\,\text{m}$ /s^2. Phản lực tại mỗi đầu bằng bao nhiêu?
\shortans{10}
\loigiai{
Do đối xứng và cân bằng: $N_1=N_2=mg/2=10$ N. Đáp án trả lời ngắn không ghi đơn vị.
}
\end{ex}

% ===== Câu 127 =====
% ID: AIQ_L10C3_FULL_1423
% Mức: VD
\begin{ex}
% ID: AIQ_L10C3_FULL_1423
% Mức: VD
Thanh đồng chất khối lượng $3\,\text{kg}$, dài $5\,\text{m}$ nằm ngang, được đỡ ở hai đầu. Tải phân bố đối xứng. Lấy g= $10\,\text{m}$ /s^2. Phản lực tại mỗi đầu bằng bao nhiêu?
\shortans{15}
\loigiai{
Do đối xứng và cân bằng: $N_1=N_2=mg/2=15$ N. Đáp án trả lời ngắn không ghi đơn vị.
}
\end{ex}

% ===== Câu 128 =====
% ID: AIQ_L10C3_FULL_1424
% Mức: VD
\begin{ex}
% ID: AIQ_L10C3_FULL_1424
% Mức: VD
Thanh đồng chất khối lượng $4\,\text{kg}$, dài $2\,\text{m}$ nằm ngang, được đỡ ở hai đầu. Tải phân bố đối xứng. Lấy g= $10\,\text{m}$ /s^2. Phản lực tại mỗi đầu bằng bao nhiêu?
\shortans{20}
\loigiai{
Do đối xứng và cân bằng: $N_1=N_2=mg/2=20$ N. Đáp án trả lời ngắn không ghi đơn vị.
}
\end{ex}

% ===== Câu 129 =====
% ID: AIQ_L10C3_FULL_1425
% Mức: VDC
\begin{ex}
% ID: AIQ_L10C3_FULL_1425
% Mức: VDC
Thanh đồng chất khối lượng $5\,\text{kg}$, dài $3\,\text{m}$ nằm ngang, được đỡ ở hai đầu. Tải phân bố đối xứng. Lấy g= $10\,\text{m}$ /s^2. Phản lực tại mỗi đầu bằng bao nhiêu?
\shortans{25}
\loigiai{
Do đối xứng và cân bằng: $N_1=N_2=mg/2=25$ N. Đáp án trả lời ngắn không ghi đơn vị.
}
\end{ex}

% ===== Câu 130 =====
% ID: AIQ_L10C3_FULL_1426
% Mức: TH
\begin{bt}
% ID: AIQ_L10C3_FULL_1426
% Mức: TH
Thanh đồng chất khối lượng $6\,\text{kg}$, dài $4\,\text{m}$ nằm ngang, được đỡ ở hai đầu. Tải phân bố đối xứng. Lấy g= $10\,\text{m}$ /s^2. Phản lực tại mỗi đầu bằng bao nhiêu?
\loigiai{
Phân tích dữ kiện, chọn định luật phù hợp. Do đối xứng và cân bằng: $N_1=N_2=mg/2=30$ N.
}
\end{bt}

% ===== Câu 131 =====
% ID: AIQ_L10C3_FULL_1427
% Mức: TH
\begin{bt}
% ID: AIQ_L10C3_FULL_1427
% Mức: TH
Thanh đồng chất khối lượng $7\,\text{kg}$, dài $5\,\text{m}$ nằm ngang, được đỡ ở hai đầu. Tải phân bố đối xứng. Lấy g= $10\,\text{m}$ /s^2. Phản lực tại mỗi đầu bằng bao nhiêu?
\loigiai{
Phân tích dữ kiện, chọn định luật phù hợp. Do đối xứng và cân bằng: $N_1=N_2=mg/2=35$ N.
}
\end{bt}

% ===== Câu 132 =====
% ID: AIQ_L10C3_FULL_1428
% Mức: VD
\begin{bt}
% ID: AIQ_L10C3_FULL_1428
% Mức: VD
Thanh đồng chất khối lượng $2\,\text{kg}$, dài $2\,\text{m}$ nằm ngang, được đỡ ở hai đầu. Tải phân bố đối xứng. Lấy g= $10\,\text{m}$ /s^2. Phản lực tại mỗi đầu bằng bao nhiêu?
\loigiai{
Phân tích dữ kiện, chọn định luật phù hợp. Do đối xứng và cân bằng: $N_1=N_2=mg/2=10$ N.
}
\end{bt}

% ===== Câu 133 =====
% ID: AIQ_L10C3_FULL_1429
% Mức: VD
\begin{bt}
% ID: AIQ_L10C3_FULL_1429
% Mức: VD
Thanh đồng chất khối lượng $3\,\text{kg}$, dài $3\,\text{m}$ nằm ngang, được đỡ ở hai đầu. Tải phân bố đối xứng. Lấy g= $10\,\text{m}$ /s^2. Phản lực tại mỗi đầu bằng bao nhiêu?
\loigiai{
Phân tích dữ kiện, chọn định luật phù hợp. Do đối xứng và cân bằng: $N_1=N_2=mg/2=15$ N.
}
\end{bt}

% ===== Câu 134 =====
% ID: AIQ_L10C3_FULL_1430
% Mức: VD
\begin{bt}
% ID: AIQ_L10C3_FULL_1430
% Mức: VD
Thanh đồng chất khối lượng $4\,\text{kg}$, dài $4\,\text{m}$ nằm ngang, được đỡ ở hai đầu. Tải phân bố đối xứng. Lấy g= $10\,\text{m}$ /s^2. Phản lực tại mỗi đầu bằng bao nhiêu?
\loigiai{
Phân tích dữ kiện, chọn định luật phù hợp. Do đối xứng và cân bằng: $N_1=N_2=mg/2=20$ N.
}
\end{bt}

% ===== Câu 135 =====
% ID: AIQ_L10C3_FULL_1431
% Mức: VDC
\begin{bt}
% ID: AIQ_L10C3_FULL_1431
% Mức: VDC
Thanh đồng chất khối lượng $5\,\text{kg}$, dài $5\,\text{m}$ nằm ngang, được đỡ ở hai đầu. Tải phân bố đối xứng. Lấy g= $10\,\text{m}$ /s^2. Phản lực tại mỗi đầu bằng bao nhiêu?
\loigiai{
Phân tích dữ kiện, chọn định luật phù hợp. Do đối xứng và cân bằng: $N_1=N_2=mg/2=25$ N.
}
\end{bt}

% ===== Câu 136 =====
% ID: AIQ_L10C3_FULL_1432
% Mức: NB
\dangbt{Vận dụng mômen lực trong dụng cụ và đời sống}
\begin{ex}
% ID: AIQ_L10C3_FULL_1432
% Mức: NB
Cùng tác dụng lực 33 $N$ vuông góc vào tay cầm, tăng khoảng cách đến trục từ $0,19\,\text{m}$ lên $0,29\,\text{m}$. Mômen tăng bao nhiêu lần?
\choice
{\True 1,53 lần}
{1,84 lần}
{1,22 lần}
{3,53 lần}
\loigiai{
Vì $M=Fd$ và $F$ không đổi nên $M_2/M_1=d_2/d_1=1,53$. Chọn A.
}
\end{ex}

% ===== Câu 137 =====
% ID: AIQ_L10C3_FULL_1433
% Mức: NB
\begin{ex}
% ID: AIQ_L10C3_FULL_1433
% Mức: NB
Cùng tác dụng lực 35 $N$ vuông góc vào tay cầm, tăng khoảng cách đến trục từ $0,21\,\text{m}$ lên $0,31\,\text{m}$. Mômen tăng bao nhiêu lần?
\choice
{1,78 lần}
{1,18 lần}
{3,48 lần}
{\True 1,48 lần}
\loigiai{
Vì $M=Fd$ và $F$ không đổi nên $M_2/M_1=d_2/d_1=1,48$. Chọn D.
}
\end{ex}

% ===== Câu 138 =====
% ID: AIQ_L10C3_FULL_1434
% Mức: NB
\begin{ex}
% ID: AIQ_L10C3_FULL_1434
% Mức: NB
Cùng tác dụng lực 37 $N$ vuông góc vào tay cầm, tăng khoảng cách đến trục từ $0,23\,\text{m}$ lên $0,33\,\text{m}$. Mômen tăng bao nhiêu lần?
\choice
{1,14 lần}
{3,43 lần}
{\True 1,43 lần}
{1,72 lần}
\loigiai{
Vì $M=Fd$ và $F$ không đổi nên $M_2/M_1=d_2/d_1=1,43$. Chọn C.
}
\end{ex}

% ===== Câu 139 =====
% ID: AIQ_L10C3_FULL_1435
% Mức: TH
\begin{ex}
% ID: AIQ_L10C3_FULL_1435
% Mức: TH
Cùng tác dụng lực 39 $N$ vuông góc vào tay cầm, tăng khoảng cách đến trục từ $0,15\,\text{m}$ lên $0,25\,\text{m}$. Mômen tăng bao nhiêu lần?
\choice
{3,67 lần}
{\True 1,67 lần}
{2 lần}
{1,34 lần}
\loigiai{
Vì $M=Fd$ và $F$ không đổi nên $M_2/M_1=d_2/d_1=1,67$. Chọn B.
}
\end{ex}

% ===== Câu 140 =====
% ID: AIQ_L10C3_FULL_1436
% Mức: TH
\begin{ex}
% ID: AIQ_L10C3_FULL_1436
% Mức: TH
Cùng tác dụng lực 33 $N$ vuông góc vào tay cầm, tăng khoảng cách đến trục từ $0,17\,\text{m}$ lên $0,27\,\text{m}$. Mômen tăng bao nhiêu lần?
\choice
{\True 1,59 lần}
{1,91 lần}
{1,27 lần}
{3,59 lần}
\loigiai{
Vì $M=Fd$ và $F$ không đổi nên $M_2/M_1=d_2/d_1=1,59$. Chọn A.
}
\end{ex}

% ===== Câu 141 =====
% ID: AIQ_L10C3_FULL_1437
% Mức: TH
\begin{ex}
% ID: AIQ_L10C3_FULL_1437
% Mức: TH
Cùng tác dụng lực 35 $N$ vuông góc vào tay cầm, tăng khoảng cách đến trục từ $0,19\,\text{m}$ lên $0,29\,\text{m}$. Mômen tăng bao nhiêu lần?
\choice
{1,84 lần}
{1,22 lần}
{3,53 lần}
{\True 1,53 lần}
\loigiai{
Vì $M=Fd$ và $F$ không đổi nên $M_2/M_1=d_2/d_1=1,53$. Chọn D.
}
\end{ex}

% ===== Câu 142 =====
% ID: AIQ_L10C3_FULL_1438
% Mức: TH
\begin{ex}
% ID: AIQ_L10C3_FULL_1438
% Mức: TH
Cùng tác dụng lực 37 $N$ vuông góc vào tay cầm, tăng khoảng cách đến trục từ $0,21\,\text{m}$ lên $0,31\,\text{m}$. Mômen tăng bao nhiêu lần?
\choice
{1,18 lần}
{3,48 lần}
{\True 1,48 lần}
{1,78 lần}
\loigiai{
Vì $M=Fd$ và $F$ không đổi nên $M_2/M_1=d_2/d_1=1,48$. Chọn C.
}
\end{ex}

% ===== Câu 143 =====
% ID: AIQ_L10C3_FULL_1439
% Mức: VD
\begin{ex}
% ID: AIQ_L10C3_FULL_1439
% Mức: VD
Cùng tác dụng lực 39 $N$ vuông góc vào tay cầm, tăng khoảng cách đến trục từ $0,23\,\text{m}$ lên $0,33\,\text{m}$. Mômen tăng bao nhiêu lần?
\choice
{3,43 lần}
{\True 1,43 lần}
{1,72 lần}
{1,14 lần}
\loigiai{
Vì $M=Fd$ và $F$ không đổi nên $M_2/M_1=d_2/d_1=1,43$. Chọn B.
}
\end{ex}

% ===== Câu 144 =====
% ID: AIQ_L10C3_FULL_1440
% Mức: VD
\begin{ex}
% ID: AIQ_L10C3_FULL_1440
% Mức: VD
Cùng tác dụng lực 33 $N$ vuông góc vào tay cầm, tăng khoảng cách đến trục từ $0,15\,\text{m}$ lên $0,25\,\text{m}$. Mômen tăng bao nhiêu lần?
\choice
{\True 1,67 lần}
{2 lần}
{1,34 lần}
{3,67 lần}
\loigiai{
Vì $M=Fd$ và $F$ không đổi nên $M_2/M_1=d_2/d_1=1,67$. Chọn A.
}
\end{ex}

% ===== Câu 145 =====
% ID: AIQ_L10C3_FULL_1441
% Mức: VDC
\begin{ex}
% ID: AIQ_L10C3_FULL_1441
% Mức: VDC
Cùng tác dụng lực 35 $N$ vuông góc vào tay cầm, tăng khoảng cách đến trục từ $0,17\,\text{m}$ lên $0,27\,\text{m}$. Mômen tăng bao nhiêu lần?
\choice
{1,91 lần}
{1,27 lần}
{3,59 lần}
{\True 1,59 lần}
\loigiai{
Vì $M=Fd$ và $F$ không đổi nên $M_2/M_1=d_2/d_1=1,59$. Chọn D.
}
\end{ex}

% ===== Câu 146 =====
% ID: AIQ_L10C3_FULL_1442
% Mức: NB
\begin{ex}
% ID: AIQ_L10C3_FULL_1442
% Mức: NB
Cùng tác dụng lực 37 $N$ vuông góc vào tay cầm, tăng khoảng cách đến trục từ $0,19\,\text{m}$ lên $0,29\,\text{m}$. Mômen tăng bao nhiêu lần?
\choiceTF
{\True Kết quả đúng là 1,53 lần.}
{Kết quả bằng 3,06 lần.}
{\True Khi hợp lực bằng 0, vật có gia tốc bằng 0.}
{Có thể bỏ qua đơn vị của đại lượng trong kết quả vật lí.}
\loigiai{
\begin{itemchoice}
\item (Đúng) Kết quả đúng là 1,53 lần. (Vì): Vì $M=Fd$ và $F$ không đổi nên $M_2/M_1=d_2/d_1=1,53$. b) Sai vì giá trị hoặc chiều nêu ra không phù hợp. c) Đúng theo định luật II Newton. d) Sai vì kết quả vật lí phải kèm đơn vị thích hợp
\item (Sai) Kết quả bằng 3,06 lần.
\item (Đúng) Khi hợp lực bằng 0, vật có gia tốc bằng 0.
\item (Sai) Có thể bỏ qua đơn vị của đại lượng trong kết quả vật lí.
\end{itemchoice}
}
\end{ex}

% ===== Câu 147 =====
% ID: AIQ_L10C3_FULL_1443
% Mức: NB
\begin{ex}
% ID: AIQ_L10C3_FULL_1443
% Mức: NB
Cùng tác dụng lực 39 $N$ vuông góc vào tay cầm, tăng khoảng cách đến trục từ $0,21\,\text{m}$ lên $0,31\,\text{m}$. Mômen tăng bao nhiêu lần?
\choiceTF
{\True Kết quả đúng là 1,48 lần.}
{Kết quả bằng 2,96 lần.}
{\True Khi hợp lực bằng 0, vật có gia tốc bằng 0.}
{Có thể bỏ qua đơn vị của đại lượng trong kết quả vật lí.}
\loigiai{
\begin{itemchoice}
\item (Đúng) Kết quả đúng là 1,48 lần. (Vì): Vì $M=Fd$ và $F$ không đổi nên $M_2/M_1=d_2/d_1=1,48$. b) Sai vì giá trị hoặc chiều nêu ra không phù hợp. c) Đúng theo định luật II Newton. d) Sai vì kết quả vật lí phải kèm đơn vị thích hợp
\item (Sai) Kết quả bằng 2,96 lần.
\item (Đúng) Khi hợp lực bằng 0, vật có gia tốc bằng 0.
\item (Sai) Có thể bỏ qua đơn vị của đại lượng trong kết quả vật lí.
\end{itemchoice}
}
\end{ex}

% ===== Câu 148 =====
% ID: AIQ_L10C3_FULL_1444
% Mức: TH
\begin{ex}
% ID: AIQ_L10C3_FULL_1444
% Mức: TH
Cùng tác dụng lực 33 $N$ vuông góc vào tay cầm, tăng khoảng cách đến trục từ $0,23\,\text{m}$ lên $0,33\,\text{m}$. Mômen tăng bao nhiêu lần?
\choiceTF
{\True Kết quả đúng là 1,43 lần.}
{Kết quả bằng 2,86 lần.}
{\True Khi hợp lực bằng 0, vật có gia tốc bằng 0.}
{Có thể bỏ qua đơn vị của đại lượng trong kết quả vật lí.}
\loigiai{
\begin{itemchoice}
\item (Đúng) Kết quả đúng là 1,43 lần. (Vì): Vì $M=Fd$ và $F$ không đổi nên $M_2/M_1=d_2/d_1=1,43$. b) Sai vì giá trị hoặc chiều nêu ra không phù hợp. c) Đúng theo định luật II Newton. d) Sai vì kết quả vật lí phải kèm đơn vị thích hợp
\item (Sai) Kết quả bằng 2,86 lần.
\item (Đúng) Khi hợp lực bằng 0, vật có gia tốc bằng 0.
\item (Sai) Có thể bỏ qua đơn vị của đại lượng trong kết quả vật lí.
\end{itemchoice}
}
\end{ex}

% ===== Câu 149 =====
% ID: AIQ_L10C3_FULL_1445
% Mức: TH
\begin{ex}
% ID: AIQ_L10C3_FULL_1445
% Mức: TH
Cùng tác dụng lực 35 $N$ vuông góc vào tay cầm, tăng khoảng cách đến trục từ $0,15\,\text{m}$ lên $0,25\,\text{m}$. Mômen tăng bao nhiêu lần?
\choiceTF
{\True Kết quả đúng là 1,67 lần.}
{Kết quả bằng 3,34 lần.}
{\True Khi hợp lực bằng 0, vật có gia tốc bằng 0.}
{Có thể bỏ qua đơn vị của đại lượng trong kết quả vật lí.}
\loigiai{
\begin{itemchoice}
\item (Đúng) Kết quả đúng là 1,67 lần. (Vì): Vì $M=Fd$ và $F$ không đổi nên $M_2/M_1=d_2/d_1=1,67$. b) Sai vì giá trị hoặc chiều nêu ra không phù hợp. c) Đúng theo định luật II Newton. d) Sai vì kết quả vật lí phải kèm đơn vị thích hợp
\item (Sai) Kết quả bằng 3,34 lần.
\item (Đúng) Khi hợp lực bằng 0, vật có gia tốc bằng 0.
\item (Sai) Có thể bỏ qua đơn vị của đại lượng trong kết quả vật lí.
\end{itemchoice}
}
\end{ex}

% ===== Câu 150 =====
% ID: AIQ_L10C3_FULL_1446
% Mức: VD
\begin{ex}
% ID: AIQ_L10C3_FULL_1446
% Mức: VD
Cùng tác dụng lực 37 $N$ vuông góc vào tay cầm, tăng khoảng cách đến trục từ $0,17\,\text{m}$ lên $0,27\,\text{m}$. Mômen tăng bao nhiêu lần?
\choiceTF
{\True Kết quả đúng là 1,59 lần.}
{Kết quả bằng 3,18 lần.}
{\True Khi hợp lực bằng 0, vật có gia tốc bằng 0.}
{Có thể bỏ qua đơn vị của đại lượng trong kết quả vật lí.}
\loigiai{
\begin{itemchoice}
\item (Đúng) Kết quả đúng là 1,59 lần. (Vì): Vì $M=Fd$ và $F$ không đổi nên $M_2/M_1=d_2/d_1=1,59$. b) Sai vì giá trị hoặc chiều nêu ra không phù hợp. c) Đúng theo định luật II Newton. d) Sai vì kết quả vật lí phải kèm đơn vị thích hợp
\item (Sai) Kết quả bằng 3,18 lần.
\item (Đúng) Khi hợp lực bằng 0, vật có gia tốc bằng 0.
\item (Sai) Có thể bỏ qua đơn vị của đại lượng trong kết quả vật lí.
\end{itemchoice}
}
\end{ex}

% ===== Câu 151 =====
% ID: AIQ_L10C3_FULL_1447
% Mức: TH
\begin{ex}
% ID: AIQ_L10C3_FULL_1447
% Mức: TH
Cùng tác dụng lực 39 $N$ vuông góc vào tay cầm, tăng khoảng cách đến trục từ $0,19\,\text{m}$ lên $0,29\,\text{m}$. Mômen tăng bao nhiêu lần?
\shortans{1,53}
\loigiai{
Vì $M=Fd$ và $F$ không đổi nên $M_2/M_1=d_2/d_1=1,53$. Đáp án trả lời ngắn không ghi đơn vị.
}
\end{ex}

% ===== Câu 152 =====
% ID: AIQ_L10C3_FULL_1448
% Mức: TH
\begin{ex}
% ID: AIQ_L10C3_FULL_1448
% Mức: TH
Cùng tác dụng lực 33 $N$ vuông góc vào tay cầm, tăng khoảng cách đến trục từ $0,21\,\text{m}$ lên $0,31\,\text{m}$. Mômen tăng bao nhiêu lần?
\shortans{1,48}
\loigiai{
Vì $M=Fd$ và $F$ không đổi nên $M_2/M_1=d_2/d_1=1,48$. Đáp án trả lời ngắn không ghi đơn vị.
}
\end{ex}

% ===== Câu 153 =====
% ID: AIQ_L10C3_FULL_1449
% Mức: VD
\begin{ex}
% ID: AIQ_L10C3_FULL_1449
% Mức: VD
Cùng tác dụng lực 35 $N$ vuông góc vào tay cầm, tăng khoảng cách đến trục từ $0,23\,\text{m}$ lên $0,33\,\text{m}$. Mômen tăng bao nhiêu lần?
\shortans{1,43}
\loigiai{
Vì $M=Fd$ và $F$ không đổi nên $M_2/M_1=d_2/d_1=1,43$. Đáp án trả lời ngắn không ghi đơn vị.
}
\end{ex}

% ===== Câu 154 =====
% ID: AIQ_L10C3_FULL_1450
% Mức: VD
\begin{ex}
% ID: AIQ_L10C3_FULL_1450
% Mức: VD
Cùng tác dụng lực 37 $N$ vuông góc vào tay cầm, tăng khoảng cách đến trục từ $0,15\,\text{m}$ lên $0,25\,\text{m}$. Mômen tăng bao nhiêu lần?
\shortans{1,67}
\loigiai{
Vì $M=Fd$ và $F$ không đổi nên $M_2/M_1=d_2/d_1=1,67$. Đáp án trả lời ngắn không ghi đơn vị.
}
\end{ex}

% ===== Câu 155 =====
% ID: AIQ_L10C3_FULL_1451
% Mức: VD
\begin{ex}
% ID: AIQ_L10C3_FULL_1451
% Mức: VD
Cùng tác dụng lực 39 $N$ vuông góc vào tay cầm, tăng khoảng cách đến trục từ $0,17\,\text{m}$ lên $0,27\,\text{m}$. Mômen tăng bao nhiêu lần?
\shortans{1,59}
\loigiai{
Vì $M=Fd$ và $F$ không đổi nên $M_2/M_1=d_2/d_1=1,59$. Đáp án trả lời ngắn không ghi đơn vị.
}
\end{ex}

% ===== Câu 156 =====
% ID: AIQ_L10C3_FULL_1452
% Mức: VDC
\begin{ex}
% ID: AIQ_L10C3_FULL_1452
% Mức: VDC
Cùng tác dụng lực 33 $N$ vuông góc vào tay cầm, tăng khoảng cách đến trục từ $0,19\,\text{m}$ lên $0,29\,\text{m}$. Mômen tăng bao nhiêu lần?
\shortans{1,53}
\loigiai{
Vì $M=Fd$ và $F$ không đổi nên $M_2/M_1=d_2/d_1=1,53$. Đáp án trả lời ngắn không ghi đơn vị.
}
\end{ex}

% ===== Câu 157 =====
% ID: AIQ_L10C3_FULL_1453
% Mức: TH
\begin{bt}
% ID: AIQ_L10C3_FULL_1453
% Mức: TH
Cùng tác dụng lực 35 $N$ vuông góc vào tay cầm, tăng khoảng cách đến trục từ $0,21\,\text{m}$ lên $0,31\,\text{m}$. Mômen tăng bao nhiêu lần?
\loigiai{
Phân tích dữ kiện, chọn định luật phù hợp. Vì $M=Fd$ và $F$ không đổi nên $M_2/M_1=d_2/d_1=1,48$.
}
\end{bt}

% ===== Câu 158 =====
% ID: AIQ_L10C3_FULL_1454
% Mức: TH
\begin{bt}
% ID: AIQ_L10C3_FULL_1454
% Mức: TH
Cùng tác dụng lực 37 $N$ vuông góc vào tay cầm, tăng khoảng cách đến trục từ $0,23\,\text{m}$ lên $0,33\,\text{m}$. Mômen tăng bao nhiêu lần?
\loigiai{
Phân tích dữ kiện, chọn định luật phù hợp. Vì $M=Fd$ và $F$ không đổi nên $M_2/M_1=d_2/d_1=1,43$.
}
\end{bt}

% ===== Câu 159 =====
% ID: AIQ_L10C3_FULL_1455
% Mức: VD
\begin{bt}
% ID: AIQ_L10C3_FULL_1455
% Mức: VD
Cùng tác dụng lực 39 $N$ vuông góc vào tay cầm, tăng khoảng cách đến trục từ $0,15\,\text{m}$ lên $0,25\,\text{m}$. Mômen tăng bao nhiêu lần?
\loigiai{
Phân tích dữ kiện, chọn định luật phù hợp. Vì $M=Fd$ và $F$ không đổi nên $M_2/M_1=d_2/d_1=1,67$.
}
\end{bt}

% ===== Câu 160 =====
% ID: AIQ_L10C3_FULL_1456
% Mức: VD
\begin{bt}
% ID: AIQ_L10C3_FULL_1456
% Mức: VD
Cùng tác dụng lực 33 $N$ vuông góc vào tay cầm, tăng khoảng cách đến trục từ $0,17\,\text{m}$ lên $0,27\,\text{m}$. Mômen tăng bao nhiêu lần?
\loigiai{
Phân tích dữ kiện, chọn định luật phù hợp. Vì $M=Fd$ và $F$ không đổi nên $M_2/M_1=d_2/d_1=1,59$.
}
\end{bt}

% ===== Câu 161 =====
% ID: AIQ_L10C3_FULL_1457
% Mức: VD
\begin{bt}
% ID: AIQ_L10C3_FULL_1457
% Mức: VD
Cùng tác dụng lực 35 $N$ vuông góc vào tay cầm, tăng khoảng cách đến trục từ $0,19\,\text{m}$ lên $0,29\,\text{m}$. Mômen tăng bao nhiêu lần?
\loigiai{
Phân tích dữ kiện, chọn định luật phù hợp. Vì $M=Fd$ và $F$ không đổi nên $M_2/M_1=d_2/d_1=1,53$.
}
\end{bt}

% ===== Câu 162 =====
% ID: AIQ_L10C3_FULL_1458
% Mức: VDC
\begin{bt}
% ID: AIQ_L10C3_FULL_1458
% Mức: VDC
Cùng tác dụng lực 37 $N$ vuông góc vào tay cầm, tăng khoảng cách đến trục từ $0,21\,\text{m}$ lên $0,31\,\text{m}$. Mômen tăng bao nhiêu lần?
\loigiai{
Phân tích dữ kiện, chọn định luật phù hợp. Vì $M=Fd$ và $F$ không đổi nên $M_2/M_1=d_2/d_1=1,48$.
}
\end{bt}
